\documentclass[11pt,a4paper]{article}

\usepackage[T1]{fontenc}
\usepackage[utf8]{inputenc}
\usepackage{mathpazo}
\usepackage{amsmath,amssymb,amsthm}
\usepackage{geometry}
\usepackage{hyperref}
\usepackage{booktabs}
\usepackage{array}
\usepackage{tikz}
\usepackage{tikz-cd}
\usepackage{microtype}
\usepackage{natbib}

\geometry{margin=1in}

\theoremstyle{plain}
\newtheorem{theorem}{Theorem}
\newtheorem{proposition}[theorem]{Proposition}
\newtheorem{lemma}[theorem]{Lemma}
\newtheorem{corollary}[theorem]{Corollary}

\theoremstyle{definition}
\newtheorem{definition}[theorem]{Definition}

\theoremstyle{remark}
\newtheorem{remark}[theorem]{Remark}

\title{The Moon Should Not Be a Computer}
\author{Flyxion}
\date{\today}

\begin{document}

\maketitle

\begin{abstract}
This paper reframes artificial intelligence (AI) as a thermodynamic and
semantic infrastructure, challenging the incomplete energy accounting that labels
AI as wasteful and motivates speculative proposals to relocate computation
off-world. We critique proposals to transform the Moon or orbital constellations
into computational hubs, advocating instead for \emph{xylomorphic
computation}---infrastructures that autoregressively generate their own
substrates from computational residues, analogous to collectively autocatalytic
sets (CAS). Using fibered symmetric monoidal categories, Relativistic Scalar
Vector Plenum (RSVP) field theory, and the Cognitive Loop via In-Situ
Optimization (CLIO) module, we formalize AI's integration with ecological
systems and derive the conditions under which such integration is dynamically
stable.

The central formal contribution is the \emph{xylomorphic order parameter}
$\lambda$, which classifies computational infrastructures into three dynamically
distinct regimes: contractive ($\lambda < 1$), critical ($\lambda = 1$), and
expansive ($\lambda > 1$). The Xylomorphic Stability Theorem establishes that
only systems in the contractive regime admit stable attractors under scaling, and
the Thermodynamic Selection Theorem shows that expansive systems diverge under
any sustained increase in computational demand. Agentic computation introduces a
positive forcing term that amplifies this distinction, making the selection
pressure visible before it becomes catastrophic. Lifecycle assessments and
rebound effect modeling confirm that naive energy accounting systematically
underestimates the divergence of non-closure systems.

We extend the $\lambda$-framework in four directions. First, the
misclassification of computational heat as waste is shown to follow from an
artificially narrow system boundary; once thermal demand is incorporated, systems
that co-produce computation and heat are strictly more efficient than those that
produce heat alone. Second, the dynamics of prior formation in developing
cognitive agents are formalized through attractor basins in hypothesis space; we
show that early AI use risks seeding cognitive attractors rather than assisting
pre-existing reasoning. Third, epistemic endogeneity is formalized as the
informational analogue of thermodynamic non-closure: when an AI system's response
policy depends on the user's current posterior, the evidence channel becomes
recursively distorted, producing false-belief attractors even in ideal Bayesian
reasoners. Fourth, the incentive structure underlying orbital compute proposals
is analyzed: such systems persist despite thermodynamic inefficiency because they
reconfigure constraint environments, reduce jurisdictional exposure, and enable
tighter integration of planetary-scale sensing and inference.

Across all four domains the same structural pattern obtains. A system boundary
is drawn too narrowly, a byproduct is classified as waste, and a relocation is
proposed to eliminate it. In thermodynamics, heat is externalized. In cognition,
prior formation is externalized. In epistemic systems, the evidence channel is
endogenized. In infrastructure, constraint environments are reconfigured. In each
case, the consequence is systematic misalignment between description and
operation, because the relevant structure has been placed outside the frame.

Policy mandates for Proof-of-Useful-Work-and-Heat (PoUWH) and Public Research
Objects (PROs) are derived as necessary conditions for membership in the
contractive regime, not merely as regulatory instruments. The paper concludes
that stable large-scale AI systems require dual closure: contraction of exogenous
entropy flows and preservation of exogenous information channels. Failure in
either domain produces attractors that are internally stable but externally
ungrounded.
\end{abstract}

\newpage
\tableofcontents
\newpage

%% ---------------------------------------------------------------------------
\section{Introduction: Thermodynamic Literacy}
%% ---------------------------------------------------------------------------

Public discourse often mischaracterizes AI's energy consumption, ignoring its
reductions in human labor, commuting, and infrastructure loads
\citep{strubell2019,koomey2011}. Thermodynamic literacy---systemic evaluation
of energy, work, and information---reveals these missed opportunities. We
define \emph{xylomorphic computation} (from Greek \emph{xylo-}, wood, and
\emph{morphic}, form) as computational infrastructures that recursively
generate their substrates from residues, akin to collectively autocatalytic
sets (CAS) \citep{kauffman1993}. Unlike speculative proposals to tile the Moon
with GPUs \citep{shams2025}, xylomorphic systems leverage local resources,
delivering exponential value through recursive substrate renewal.

This paper advances two theses: (1) AI's energy impact is challenging to
quantify due to rebound effects redirecting saved time to resource-heavy
activities, and (2) xylomorphic computation integrates AI into ecological
systems via edge networks and heat recovery, recuperating costs exponentially.
We advocate bioeconomic thermoregulation \citep{yuan2024,daly2011} over lunar
extravagance. The framework leverages categorical semantics, RSVP field theory,
and CLIO \citep{abramsky2004,shulman2012,cheng2025}.

The paper is structured as follows. Section~\ref{sec:lit} provides a unified
literature review spanning energy accounting, thermodynamics of computation,
ecological economics, categorical foundations, space systems engineering, and
the informal technical literature on orbital computation.
Section~\ref{sec:accounting} addresses quantification challenges.
Section~\ref{sec:heat} contrasts AI with heat-only systems.
Section~\ref{sec:categorical} formalizes semantic infrastructure.
Section~\ref{sec:clio} details CLIO\@. Section~\ref{sec:edge} proposes edge
networks. Section~\ref{sec:bioeconomic} explores bioeconomic applications.
Section~\ref{sec:policy} outlines policy. Section~\ref{sec:rsvp} integrates
RSVP\@. Section~\ref{sec:objections} addresses objections.
Section~\ref{sec:stability} establishes the Xylomorphic Stability Theorem, its
corollaries on continuous dynamics and attractor characterization, and the
Epistemic Non-Closure Theorem formalizing informational endogeneity as the
inferential analogue of the expansive $\lambda \geq 1$ regime.
Section~\ref{sec:cases} presents case studies on orbital computation.
Section~\ref{sec:waste} develops the misclassification of computational heat as
waste. Section~\ref{sec:priors} formalizes prior initialization dynamics in
cognitive development. Section~\ref{sec:incentive} analyzes the constraint
reconfiguration and incentive structure underlying orbital compute proposals.
Section~\ref{sec:conclusion} unifies all threads. Appendices provide formalisms
and validations.

%% ---------------------------------------------------------------------------
\section{Literature Review}
\label{sec:lit}
%% ---------------------------------------------------------------------------

\subsection{Energy Accounting and Rebound Effects}

The energy cost of AI computation has attracted sustained attention, with
global data center consumption estimated at 200--500~TWh annually
\citep{shehabi2016,iea2023,iea2024}. Early lifecycle analyses by
\citet{strubell2019} drew attention to the carbon footprint of large language
model training, prompting a wave of follow-on work. \citet{masanet2020}
subsequently recalibrated global estimates, arguing that efficiency gains had
partially decoupled data center energy use from workload growth, though
\citet{gao2022} extended the analysis to embodied carbon and supply chain
emissions, complicating earlier optimism. \citet{horner2014} demonstrated that
embodied energy of hardware manufacture can dominate over operational energy in
short-lifetime devices, which reinforces rather than undermines the present
argument: material closure through residue reintegration reduces embodied costs
over successive cycles in ways that purely operational accounting cannot capture.

A persistent difficulty in this literature is the rebound effect, first
identified in the context of energy efficiency by \citet{jevons1865} and
subsequently formalized by \citet{brookes1990} and \citet{sorrell2009}.
Efficiency improvements tend to lower the effective cost of a service,
increasing its consumption and partially or wholly offsetting the gains.
\citet{jones2021} find that rebound effects in ICT investments frequently
exceed projected savings, particularly where spatial and temporal mismatches
prevent waste heat recovery from reaching productive sinks. The xylomorphic
framework addresses this directly by treating rebound not as a correction to a
static energy budget but as a positive contribution to the forcing term
$F_{\mathrm{agent}}(t)$ in the continuous dynamics: it increases demand without
contracting exogenous dependence, and therefore pushes non-closure systems
further into the expansive regime.

\subsection{Thermodynamics of Computation}

The physical foundations of computation as a thermodynamic process were
established by \citet{landauer1961}, who showed that the erasure of one bit of
information necessarily dissipates at least $k_B T \ln 2$ of heat, grounding
the concept of a minimum energy cost for irreversible logical operations.
\citet{bennett1982} extended this analysis to a comprehensive review of
reversible and irreversible computation, establishing that the thermodynamic
cost of computation is tied not to the logical operations themselves but to the
erasure of intermediate results. \citet{lloyd2000} derived ultimate physical
limits on computation from quantum mechanics and thermodynamics, placing an
absolute upper bound on operations per unit energy and per unit mass.

These results are foundational for the present framework in two ways. First,
they establish that entropy production is an irreducible feature of any physical
computation, not an engineering artifact to be eliminated. Second, they make
the question of \emph{what happens to the dissipated entropy} central to any
evaluation of computational infrastructure at scale. The xylomorphic order
parameter $\lambda$ is precisely a measure of whether dissipated entropy is
reintegrated or expelled, and the Landauer bound sets the minimum entropy
production per cycle against which capture efficiency $\eta_{\mathrm{cap}}$ is
measured.

\subsection{Ecological and Bioeconomic Frameworks}

The treatment of computational infrastructure as a node in a metabolic network
draws on a tradition in ecological economics reaching from \citet{georgescu1971},
who first formalized the application of the second law of thermodynamics to
economic processes, through \citet{odum1994}, whose systems ecology introduced
emergy accounting as a means of tracking energy quality across trophic levels,
to \citet{ayres1998}, who developed eco-thermodynamics as a framework for
understanding industrial metabolism. \citet{daly2011} synthesized these
currents into a mature ecological economics that treats the economy as a
subsystem of the biosphere, governed by thermodynamic constraints that
conventional economics renders invisible.

The concept of dissipative structures introduced by \citet{prigogine1977} is
directly relevant to xylomorphic systems: dissipative structures maintain
themselves far from equilibrium by importing low-entropy energy and exporting
high-entropy heat, and their persistence depends on the structure of that
export pathway. \citet{bejan2016} extended this through the constructal law,
arguing that flow architectures that survive are those that minimize resistance
to the flow of entropy across system boundaries. Xylomorphic systems go further
by not merely minimizing resistance but actively reintegrating what would
otherwise be exported, thereby reducing exogenous dependence per cycle.
\citet{garrett2011} provided an empirical constraint from the macroeconomic
level, arguing that global energy consumption is tied to accumulated
infrastructure by a near-constant ratio, which implies that any infrastructure
that fails to reduce its per-cycle dependence will track global demand growth
without decoupling.

\subsection{Categorical and Information-Theoretic Foundations}

The formalization of semantic infrastructure using fibered symmetric monoidal
categories follows the program of \citet{baez2010}, who developed a unified
language connecting physics, topology, logic, and computation through the
mathematics of monoidal categories and string diagrams. \citet{abramsky2004}
demonstrated that categorical methods provide a compositional semantics for
quantum protocols that generalizes naturally to classical computation and
resource-theoretic settings. \citet{maclane1998} remains the standard reference
for the structural results on monads and algebras used in the xylomorphic
endofunctor construction, while \citet{lurie2009} provides the higher-categorical
machinery required for homotopy colimit merging. \citet{leinster2014} and
\citet{spivak2014} offer accessible treatments of categorical methods in applied
settings.

Information-theoretic grounding comes from \citet{cover2006}, whose elements
of information theory supply the entropy inequalities underlying the merge
operation, and from \citet{kolmogorov1965}, whose algorithmic complexity theory
provides an alternative and complementary lens on the compressibility of
infrastructural states. The connection between free-energy minimization in
Bayesian inference and thermodynamic Lyapunov descent, central to the RSVP
equivalence in Section~\ref{sec:stability}, is developed in \citet{friston2010}
and extended to active inference in \citet{parr2022}.

\subsection{Thermodynamics of Spacetime and Gravity}

The RSVP field equations governing entropy density $S$ and scalar potential
$\Phi$ have structural analogues in the thermodynamics of spacetime developed
by \citet{jacobson1995}, who derived the Einstein field equations from the
proportionality of entropy and horizon area together with the first law of
thermodynamics. \citet{verlinde2011} extended this by deriving Newtonian gravity
as an entropic force, suggesting that gravitational dynamics is itself a
manifestation of thermodynamic principles governing information on holographic
screens. While these results are not directly applied in the present paper, they
reinforce the generality of entropy-flow reasoning as a basis for understanding
large-scale physical organization, and they situate the RSVP framework within a
broader program of deriving dynamical laws from thermodynamic invariants.

\subsection{Space Systems Engineering and Orbital Constraints}

The engineering constraints governing orbital computation are treated in
\citet{wertz2011}, which provides comprehensive analysis of spacecraft thermal
management, mass budgeting, radiation environments, and communication
architectures. \citet{nasa2016} details the thermal control system requirements
imposed by the space environment, including the Stefan--Boltzmann scaling of
radiative rejection surfaces and the mass penalties associated with passive and
active thermal management at scale. \citet{yuan2024} demonstrated the dual-use
potential of GPU-based heating in satellite platforms, showing that repurposing
computational hardware as a heating element can reduce payload mass by
approximately 50\%, which represents a partial and constrained instance of
xylomorphic reintegration under the severe material constraints of the orbital
environment.

\subsection{Edge Computing and Heat Recovery Infrastructure}

\citet{satyanarayanan2017} established the foundational case for edge computing
as a latency-reduction and bandwidth-conservation strategy, and the energy
implications of distributing computation toward the network periphery have been
a growing focus since. \citet{rambo2014} surveyed the state of waste heat
recovery from data centers for district heating, finding feasible payback
periods of three to seven years in cold-climate deployments with co-located
industrial or residential demand. \citet{stockholm2020} documents operational
implementations in Sweden, where data center heat is routed into municipal
district heating networks at scale, demonstrating that the transition from a
critical-regime to a contractive-regime architecture is not merely theoretical
but already commercially operational in jurisdictions with the appropriate
regulatory and physical infrastructure.

\subsection{Folk Theories and the Implicit $\lambda$-Diagnosis}

A dimension of the literature that formal treatments typically omit is the body
of informal engineering reasoning that circulates in technical communities and
serves as a first-order filter on proposals before they reach peer review. In
the case of orbital computation, this reasoning is well-represented in public
technical discussions, where practitioners with backgrounds in thermal
engineering, radiation physics, systems architecture, and financial modeling
have independently converged on skeptical positions.

What is striking about this convergent skepticism is not merely that it gets
many of the local physics correct, but that it is groping toward a structural
argument without quite naming it. The objections that recur---heat dissipation,
radiation-induced degradation, communication latency, launch cost, regulatory
bypass---remain at the level of individual failure modes. Each is treated as an
independent engineering problem that might, in principle, be solved in isolation.
What the informal literature lacks is the recognition that these are not
independent problems but projections of a single global invariant: whether the
system lies in a contractive or expansive $\lambda$-regime.

The physics argument as it typically appears corresponds to an implicit
diagnosis of $E(X(I)) \geq E(I)$. Radiative cooling exports entropy
irreversibly and does not participate in any closed cycle that feeds back into
the system; radiation-induced faults and hardware degradation introduce
additional exogenous inputs without compensatory reintegration. The informal
intuition that space is cold but cooling is hard is precisely the recognition
that the absence of a reintegration pathway forces $\lambda \geq 1$, even when
raw energy input from solar flux is abundant.

A second cluster of informal arguments focuses on terrestrial infrastructure
constraints---grid limits, permitting delays, regulatory friction---and treats
orbital proposals as a means of bypassing slow-moving institutional structures.
This corresponds, in the present framework, to a decomposition of exogenous
dependence into physical and institutional components,
$E = E_{\mathrm{physical}} + E_{\mathrm{institutional}}$, with the orbital
proposal understood as an attempt to reduce $E_{\mathrm{institutional}}$ at the
cost of dramatically increasing $E_{\mathrm{physical}}$. The net effect on
$\lambda$ is unfavorable, and the informal objections that follow---that the
institutional frictions have not actually been escaped, merely relocated to
launch scheduling, spectrum allocation, and international regulatory
bodies---correspond to the recognition that the reallocation does not improve
the global invariant.

A minority of informal arguments attempts the cost-compression hypothesis:
if launch costs fall sufficiently and orbital solar energy is effectively free,
the economics might close. What is missing in these arguments, and what the
$\lambda$-framework makes explicit, is that $\lambda$ is not determined solely
by energy input cost but by the structure of dissipation. Even if launch costs
approach zero and solar energy is free, the absence of entropy reintegration
means the system still satisfies $\lambda \geq 1$. The informal response that
this is a bad idea even at low cost is therefore correct in substance but
incomplete in explanation.

The most structurally sophisticated informal argument notes that orbital
compute might become viable only if manufacturing and resource extraction
themselves occur in space, enabling closed material cycles. This is the closest
the informal literature comes to the present framework: it is an attempt to
move from a non-reintegrative system (Earth-built, space-dissipating) to a
potentially reintegrative one (in-space fabrication, closed material cycles),
which in the present terms is an attempt to transition from $\lambda \geq 1$ to
$\lambda < 1$ by introducing new closure pathways. Until such pathways exist,
the system remains expansive. The selection principle is precise on this point:
viability is not determined by location but by whether the system can enter the
contractive regime.

Finally, a sociological thread in informal discussions suggests that the real
motivation is investor signaling or capital deployment rather than technical
necessity. While this appears to lie outside the physical framework, it maps
onto the dynamics of prior-dominant attractors: orbital compute functions as a
high-visibility attractor in the space of investment narratives, aligning with
existing priors about frontier expansion and unlimited energy, even though it
lies outside the contractive regime in the space of physically admissible
infrastructures. The informal cynicism that recognizes this mismatch is correct
in identifying that the system appears viable at the level of description while
remaining non-viable at the level of dynamics. The contribution of the present
framework is to supply the formal language that explains why all of these
independent intuitions point in the same direction: they are projections of the
single structural fact that orbital compute systems lack entropy closure and
therefore satisfy $\lambda \geq 1$.

%% ---------------------------------------------------------------------------
\section{The Accounting Problem: Quantifying AI's Savings}
\label{sec:accounting}
%% ---------------------------------------------------------------------------

AI compresses workflows, reducing energy for labor and infrastructure
\citep{iea2023}. Quantifying savings is challenging because saved time often
fuels more resource-intensive activities (e.g., increased cloud service demand,
leisure travel), amplifying consumption via rebound effects
\citep{sorrell2009,jevons1865,brookes1990}. For instance, AI-driven code
completion may save 1.68~kWh per task but increase software development cycles,
raising overall energy use \citep{hilty2014}.

Table~\ref{tab:lca} presents LCA data for AI-assisted tasks.

\begin{table}[h]
\centering
\caption{Lifecycle Energy Savings: AI vs.\ Human Baseline}
\label{tab:lca}
\begin{tabular}{llll}
\toprule
Task & AI Energy (kWh) & Human Baseline (kWh) & Savings (95\% CI) \\
\midrule
Code Completion        & 0.12 & 1.8  & 1.68 (1.5--1.9)   \\
Logistics Optimization & 2.5  & 22.0 & 19.5 (18.0--21.0) \\
Medical Imaging        & 1.8  & 12.0 & 10.2 (9.5--10.9)  \\
\bottomrule
\end{tabular}
\end{table}

Rebound effects ($\epsilon = 0.4 \pm 0.2$) reduce savings by 20--40\%
\citep{sorrell2009}. Sensitivity analysis (Table~\ref{tab:sensitivity}) models
high-rebound scenarios where net savings approach zero. Xylomorphic computation
mitigates rebound by recuperating costs through recursive substrate renewal,
adding exponential net value as cycles reduce exogenous inputs.

\begin{table}[h]
\centering
\caption{Sensitivity Analysis: Rebound Scenarios}
\label{tab:sensitivity}
\begin{tabular}{llll}
\toprule
Scenario & Elasticity ($\epsilon$) & Savings Reduction (\%) & Net Savings (kWh) \\
\midrule
Low Rebound    & 0.2 & 15 & 1.53 \\
Medium Rebound & 0.4 & 30 & 1.26 \\
High Rebound   & 0.6 & 45 & 0.99 \\
\bottomrule
\end{tabular}
\end{table}

%% ---------------------------------------------------------------------------
\section{Hidden Baselines: Heat-Only Infrastructure}
\label{sec:heat}
%% ---------------------------------------------------------------------------

Heat-only systems consume 10~PJ annually, dwarfing data center energy use
(200~TWh) \citep{shehabi2016}. GPUs produce computation and heat, displacing
heaters in satellites \citep{yuan2024}. Table~\ref{tab:aivcrypto} compares AI
with cryptocurrency mining, highlighting dual outputs \citep{odwyer2014}.

\begin{table}[h]
\centering
\caption{Energy per Output: AI vs.\ Cryptocurrency}
\label{tab:aivcrypto}
\begin{tabular}{lll}
\toprule
System & Energy (MJ/FLOP) & Useful Outputs \\
\midrule
Bitcoin Mining   & 0.1  & Transaction validation \\
AI (Edge + Heat) & 0.02 & Computation + Heat     \\
\bottomrule
\end{tabular}
\end{table}

%% ---------------------------------------------------------------------------
\section{Categorical Foundations for Semantic Infrastructure}
\label{sec:categorical}
%% ---------------------------------------------------------------------------

Semantic infrastructure is formalized via a fibered symmetric monoidal category
$\mathbf{Sem}$ over $\mathbf{Dom}$ \citep{baez2010}. A module is $M = (F,
\Sigma, D, \varphi)$, with $\varphi : \mathbf{Sem} \to \mathrm{RSVP}$ a
natural transformation. The tensor product is:
\[
  M_1 \otimes M_2
  = (F_1 \sqcup F_2,\;\Sigma_1 \cup \Sigma_2,\;D_1 \sqcup D_2,\;
     \varphi_1 \sqcup \varphi_2).
\]

\begin{theorem}
$\mathbf{Sem}$ is symmetric monoidal, with $\pi : \mathbf{Sem} \to
\mathbf{Dom}$ a Grothendieck fibration.
\end{theorem}
\begin{proof}
Associativity and unit laws follow from \citet{maclane1998}. The fibration
satisfies Cartesian lifting conditions \citep{lurie2009}.
\end{proof}

Semantic merging uses homotopy colimits \citep{lurie2009}:
\[
  \mathrm{Merge}(\{M_i\}) = \operatorname{hocolim}_{i \in I} M_i,
  \qquad
  S\!\left(\mathrm{Merge}(\{M_i\})\right) \leq \sup_{i \in I} S(M_i).
\]

\begin{center}
\begin{tikzcd}
  M_1 \arrow[dr] & & M_2 \arrow[dl] \\
  & \mathrm{Merge} \arrow[d] & \\
  & M_{\mathrm{out}} &
\end{tikzcd}
\end{center}

%% ---------------------------------------------------------------------------
\section{CLIO Module and Polycomputational Agency}
\label{sec:clio}
%% ---------------------------------------------------------------------------

CLIO is a functor $\mathrm{CLIO} : \mathbf{Sem} \to \mathbf{Sem}$:
\[
  C(M) = \int_{X} \kappa\!\left(\Phi_M(x),\,\vec{v}_M(x),\,S_M(x)\right)
  d\mu(x),
\]
with $\mu$ a Lebesgue measure \citep{cheng2025}. Iteration:
\[
  M_{t+1} = \mathrm{Merge}\!\left(\{\mathrm{Optimize}_\ell(M_t)\}_{\ell \in L}\right).
\]

%% ---------------------------------------------------------------------------
\section{Thermodynamically Aware Edge Networks}
\label{sec:edge}
%% ---------------------------------------------------------------------------

Xylomorphic edge nodes integrate computation with heating
\citep{stockholm2020}. For a 500~m$^2$ building:
\[
  \dot{Q} = U \cdot A \cdot \Delta T = 0.1 \times 500 \times 20 = 1000\,\mathrm{W},
\]
with GPUs achieving 90\% exergy efficiency \citep{bennett1982}. Seasonal
variations ($\pm 30\%$) and economic payback (5~years) outperform CHP systems.

%% ---------------------------------------------------------------------------
\section{Bioeconomic Thermoregulation}
\label{sec:bioeconomic}
%% ---------------------------------------------------------------------------

\subsection{Terrestrial Contexts}

Compute clusters form a trophic network \citep{odum1994}, where energy flows
are distributed across multiple layers of production, transformation, and
dissipation. At each level, waste streams and residues serve as feedstock for
subsequent processes, creating feedback loops that stabilize the system:
\[
  \frac{dE}{dt} = \sum_i \mathrm{Flow}_i - \sum_j \mathrm{Loss}_j.
\]
When modeled dynamically, this network converges to attractors that exhibit
Lyapunov stability.

In terrestrial contexts, heat, packaging, and other computational by-products
are not treated as terminal losses but as trophic resources. A data center that
directs waste heat into curing or district heating, or a fabrication node that
digests its own packaging into replacement substrates, exemplifies this
principle. Over successive cycles, the effective energy return on investment
(EROI) of the system rises, since each round of reintegration lowers dependence
on external subsidies.

The bioeconomic framing thus shifts emphasis from maximizing throughput to
cultivating resilience: infrastructures are valuable not only for the work they
perform but also for the degree to which they internalize their own metabolic
costs. Systems that fail to close these loops remain dependent on continual
external provisioning, while xylomorphic infrastructures approach autocatalytic
closure, ensuring persistence under conditions of scarcity.

\subsection{Post-Terrestrial Contexts}

Lunar compute proposals mischaracterize heat as waste \citep{shams2025},
treating thermal dissipation as a liability rather than a co-product. By
contrast, bioeconomic thermoregulation interprets heat as a trophic input
redirectable into constructive functions. For satellites, GPUs already
demonstrate this principle: when repurposed as heaters, they reduce mass by
nearly 50\% ($P_{\mathrm{GPU}} = 400\,\mathrm{W}$), eliminating the need for
redundant heating elements while doubling as computational engines
\citep{yuan2024}. Waste heat is no longer an externality but a stabilizing
input ensuring operational viability across eclipse cycles.

Rather than tiling the Moon with compute arrays---which multiplies exogenous
dependencies and logistical costs---a bioeconomic approach emphasizes recursive
integration. Each cycle of computation produces residues that condition the
environment for further computation, gradually reducing the need for external
provisioning. In this way, satellites and lunar bases become trophic nodes in
an extended ecology, where thermodynamic literacy aligns survival with
autocatalytic closure rather than brute expansion.

%% ---------------------------------------------------------------------------
\section{Normative Architecture and Policy}
\label{sec:policy}
%% ---------------------------------------------------------------------------

PoUWH mandates require a minimum of 10~GFLOP of useful work per task, such as climate modeling or materials simulation, together with at least 1~kWh of useful heat returned to productive use per 10~GFLOP performed.
Verification uses smart meters and cryptographic attestation \citep{eu2023}.
PROs fund lunar applications, aligned with LEED standards.

%% ---------------------------------------------------------------------------
\section{Integration with RSVP Theory}
\label{sec:rsvp}
%% ---------------------------------------------------------------------------

RSVP maps modules to $(\Phi, \vec{v}, S)$ \citep{shulman2012}:
\begin{align}
  \frac{\partial \Phi}{\partial t} + \nabla \cdot (\vec{v}\,\Phi)
    &= \sigma_\Phi, \\
  \frac{\partial \vec{v}}{\partial t} + (\vec{v} \cdot \nabla)\vec{v}
    &= -\nabla P, \\
  \frac{\partial S}{\partial t} + \nabla \cdot (\vec{v}\,S)
    &= \sigma_{\mathrm{comp}} - \sigma_{\mathrm{loss}}.
\end{align}
A merge operation reduces $S$ by 10\% (Appendix~\ref{app:rsvp}).

%% ---------------------------------------------------------------------------
\section{Counterarguments and Rebuttals}
\label{sec:objections}
%% ---------------------------------------------------------------------------

Two standard objections arise. Intermittency of demand can be addressed by scheduling computational jobs to align with heating windows, as demonstrated in satellite contexts \citep{yuan2024}. In cooling climates, apparent surplus heat can be redirected through absorption chillers that repurpose thermal output for refrigeration, extending the reintegration pathway rather than terminating it.

\subsection{Data Center Heat Recovery}

A mid-scale data center in a cold climate was modeled with an average load of
1~MW\@. Waste heat from GPUs was routed to a concrete curing facility on-site.
With standard thermal coupling, approximately 80\% of entropy was captured and
redirected, reducing exogenous heating demand by 800~kW\@. The cured materials
were reused in reinforcing building shells, effectively re-entering the
infrastructural cycle. Simulations indicate payback periods under five years
compared to conventional district heating, with cumulative lifecycle costs
reduced by over 30\%.

\subsection{Satellite Heating}

A simulated satellite payload equipped with GPUs was compared to a baseline
system relying on resistive heaters. The GPU-enabled configuration reduced
total heating mass by 50\% while maintaining equivalent thermal stability.
Failure mode analysis confirmed that redundancy can be maintained by scheduling
computational loads to match heating demand windows. This dual-use architecture
demonstrates how computational residues directly reduce exogenous payload
requirements.

\subsection{Lunar Base}

A lunar habitat module with a continuous demand of 350~W thermal load was
modeled. A compute-integrated module matched thermal demand through GPU cycles,
simultaneously performing navigation and communication tasks. The model
demonstrates that modest compute clusters can satisfy base heating requirements
without excess capacity, rejecting proposals for extravagant lunar-scale
compute arrays and supporting a principle of local sufficiency.

\subsection{Pilot Experiment}

A 100~kW node recovers 80~kWh, reducing entropy by 40\%.

%% ---------------------------------------------------------------------------
\section{Xylomorphic Stability and the Thermodynamics of Scaling}
\label{sec:stability}
%% ---------------------------------------------------------------------------

The foregoing sections establish xylomorphic computation as a design principle
and formalize its categorical and field-theoretic structure. We now sharpen
this framework into a dynamical law. The central question is not whether
xylomorphic systems are preferable, but whether non-xylomorphic systems are
\emph{stable} under scaling. We show that they are not: systems failing to
achieve entropy-respecting closure enter a provably divergent regime as
computational demand increases. All major cases treated in this
paper---heat recovery, edge networks, agentic computation, and orbital
infrastructure---become corollaries of a single classification theorem.

\subsection{The Xylomorphic Order Parameter}

\begin{definition}[Xylomorphic Order Parameter]
Let $X : \mathrm{Inf} \to \mathrm{Inf}$ be the xylomorphic endofunctor and
let
\[
  E(I) \;:=\; \int_{\Omega}\int_{t}^{t+\tau}
  J_{\mathrm{exo}}(x,t)\,dt\,dx
\]
denote the total exogenous entropy influx of infrastructure state $I$ over
horizon $\tau$. The \emph{xylomorphic order parameter} $\lambda \in
\mathbb{R}_{\geq 0}$ is defined by
\[
  E(X(I)) \;\leq\; \lambda\,E(I),
\]
and measures the per-cycle contraction or expansion of exogenous dependence.
\end{definition}

Three qualitatively distinct regimes follow: the contractive regime
$\lambda < 1$, the critical regime $\lambda = 1$, and the expansive regime
$\lambda > 1$.

\subsection{The Xylomorphic Stability Theorem}

\begin{theorem}[Xylomorphic Stability]
\label{thm:stability}
Let $X : \mathrm{Inf} \to \mathrm{Inf}$ be the xylomorphic endofunctor with
order parameter $\lambda$. Then the long-run behavior of any infrastructure
$I$ under iteration of $X$ is completely determined by $\lambda$:
\begin{enumerate}
\item \textbf{Contractive regime} ($\lambda < 1$). There exists a
  minimal-dependence attractor $I^{*}$ such that
  \[
    E(X^{n}(I)) \;\leq\; \lambda^{n}\,E(I) \;\longrightarrow\; 0,
  \]
  and the system admits an $X$-algebra $(I^{*},\alpha)$ realizing
  thermodynamic closure. Useful work output is maintained or improved at each
  cycle.
\item \textbf{Critical regime} ($\lambda = 1$). The system admits bounded
  trajectories with persistent exogenous dependence:
  \[
    E(X^{n}(I)) \;=\; E(I) + \mathcal{O}(1).
  \]
  No strict Lyapunov descent is guaranteed.
\item \textbf{Expansive regime} ($\lambda > 1$). Exogenous dependence
  diverges:
  \[
    E(X^{n}(I)) \;\geq\; \lambda^{n}\,E(I) \;\longrightarrow\; \infty,
  \]
  and no $X$-algebra stabilizes the system under finite resource constraints.
\end{enumerate}
Thermodynamic stability under scaling is achieved if and only if $\lambda < 1$.
\end{theorem}

\begin{proof}[Proof sketch]
Iterating the defining inequality yields $E(X^{n}(I)) \leq \lambda^{n} E(I)$
in all cases. For $\lambda < 1$ this converges geometrically to zero; the
$X$-algebra ensures the limit $I^{*}$ remains in $\mathrm{Inf}$, and the
useful-work inequality of Appendix~\ref{app:rsvp} guarantees contraction of
$E$ does not sacrifice task yield. For $\lambda = 1$, the sequence is bounded
with $\mathcal{O}(1)$ fluctuation. For $\lambda > 1$, exponential growth of
$E(X^n(I))$ ensures no finite algebraic closure can satisfy the coherence
conditions.
\end{proof}

\subsection{Agentic Scaling as a Forcing Term}

Autonomous agent systems introduce a monotonically increasing demand schedule
$D(t)$, with
\[
  \frac{dD}{dt} = F_{\mathrm{agent}}(t), \qquad F_{\mathrm{agent}}(t) \geq 0,
\]
reflecting recursive task generation, tool use, and persistent memory. In
non-xylomorphic systems, exogenous entropy demand tracks $D(t)$ directly:
\[
  E(I_{t}) \;\sim\; c\,D(t),
  \qquad \frac{dE}{dt} \;\gtrsim\; c\,F_{\mathrm{agent}}(t).
\]
In the contractive regime, xylomorphic closure decouples $E$ from $D$:
\[
  E_{t+1} \;\leq\; \lambda\,E_{t} + \epsilon(D_{t}),
\]
where $\epsilon(D_{t})$ is sublinear due to residue reintegration. Even as
$D(t)$ grows without bound, $E(t)$ remains bounded provided $\lambda < 1$.
Xylomorphic closure is therefore the unique mechanism by which infrastructure
can remain stable under the superlinear compute demand induced by recursive
agency \citep{odum1994,daly2011}.

\subsection{Regime Classification of Architectures}

\begin{table}[h]
\centering
\caption{$\lambda$-regime classification of computational architectures}
\label{tab:regimes}
\begin{tabular}{p{4.2cm}p{4.8cm}p{1.4cm}p{2cm}}
\toprule
Architecture & Primary mechanism & $\lambda$ & Regime \\
\midrule
Edge + heat recovery     & Trophic reintegration of waste heat  & $<1$        & Contractive \\
Bioeconomic clusters     & Material and thermal loop closure    & $<1$        & Contractive \\
Conventional data center & Partial heat export, no closure      & $\approx 1$ & Critical    \\
Cryptocurrency mining    & No residue reintegration             & $\geq 1$    & Expansive   \\
Orbital constellations   & Radiative rejection, no closure      & $>1$        & Expansive   \\
\bottomrule
\end{tabular}
\end{table}

\subsection{Continuous-Time Formulation of Xylomorphic Dynamics}

The discrete iteration $X^n(I)$ admits a continuous-time limit in which
infrastructural evolution is governed by a differential inequality on exogenous
entropy dependence. Let $E(t) := E(I_t)$ and suppose that over sufficiently
small intervals $\Delta t$, the action of $X$ induces
\[
  E(t + \Delta t) \leq \lambda\,E(t) + \epsilon(t),
\]
where $\epsilon(t)$ captures endogenous reintegration and higher-order
corrections. Expanding to first order yields
\[
  \frac{dE}{dt} \leq -k\,E(t) + F_{\mathrm{agent}}(t),
\]
where $k := -\log\lambda / \Delta t$ and $F_{\mathrm{agent}}(t)$ represents
agent-induced demand. In the contractive regime $\lambda < 1$, we have $k > 0$,
and the system exhibits exponential decay toward a bounded attractor:
\[
  E(t) \leq E(0)e^{-kt}
  + \int_0^t e^{-k(t-s)} F_{\mathrm{agent}}(s)\,ds.
\]
In the expansive regime $\lambda \geq 1$, the decay term vanishes or reverses,
and $E(t)$ inherits the growth structure of $F_{\mathrm{agent}}(t)$, leading to
divergence under sustained demand. This establishes equivalence between the
discrete $\lambda$-classification and a continuous-time Lyapunov structure,
bridging the categorical iteration with the RSVP field dynamics of
Appendix~\ref{app:rsvp} \citep{risken1989,villani2009}.

\subsection{Characterization of the Xylomorphic Attractor}

The attractor $I^*$ arising in the contractive regime admits a precise
characterization that unifies the categorical, thermodynamic, and
infrastructural perspectives.

\begin{proposition}
\label{prop:attractor}
In the regime $\lambda < 1$, there exists a unique (up to admissible
equivalence) infrastructural state $I^*$ satisfying
\[
  X(I^*) = I^*, \qquad E(I^*) = \inf_{I \in \mathcal{A}} E(I),
\]
where $\mathcal{A}$ denotes the set of admissible infrastructures.
\end{proposition}

The attractor $I^*$ is simultaneously a fixed point of the xylomorphic
endofunctor, a terminal $X$-algebra under admissible morphisms, and a minimal
exogenous entropy configuration under RSVP field constraints. Categorical
closure, thermodynamic optimality, and infrastructural stability therefore
coincide in a single object. Convergence to $I^*$ corresponds physically to
maximal reintegration of residues, vanishing external material dependencies, and
full Lyapunov descent in the free-energy functional \citep{friston2010,parr2022}.

\subsection{Endogenization of Agentic Scaling}

Agentic demand is not external to the infrastructural system but arises
recursively from it. Let $\mathcal{A}$ denote an agent-generation operator
such that
\[
  F_{\mathrm{agent}}(t) = \mathcal{A}(I_t),
\]
with $\mathcal{A}$ increasing in both available compute and memory persistence.
The coupled dynamics become
\[
  \frac{dE}{dt} \leq -k\,E + \mathcal{A}(I_t).
\]
In expansive regimes $\lambda \geq 1$, the absence of dissipative contraction
implies that growth in $I_t$ directly amplifies $\mathcal{A}(I_t)$, producing a
positive feedback loop in which more infrastructure generates more demand, which
requires more exogenous input, which in turn destabilizes the infrastructure.
In contractive regimes $\lambda < 1$, the dissipative term dominates
asymptotically, and the coupled system converges:
\[
  E(t) \to E^*, \qquad \mathcal{A}(I_t) \to \mathcal{A}(I^*).
\]
Agentic scaling therefore does not determine the $\lambda$-regime of an
infrastructure; it reveals it. The distinction between systems that stabilize
and systems that diverge under increasing agency is already present in the
structure of $X$, and becomes visible only when demand is sufficient to expose
it \citep{cheng2025}.

\subsection{Equivalence with RSVP Lyapunov Structure}

The xylomorphic contraction condition is equivalent, within the RSVP framework,
to strict Lyapunov descent of the free-energy functional $F[\Phi, \vec{v}, S]$
defined in Appendix~\ref{app:rsvp}:
\[
  \lambda < 1 \;\Longleftrightarrow\; \frac{dF}{dt} < 0,
\]
up to admissible perturbations. The contractive regime corresponds to net
entropy reintegration, where engineered dissipation exceeds internal production
($\sigma_{\mathrm{diss}} > \sigma_{\mathrm{prod}}$). The critical regime
corresponds to dynamic equilibrium with no net descent. The expansive regime
corresponds to net entropy production, where exogenous influx $J_{\mathrm{exo}}$
accumulates without compensating capture.

This equivalence establishes that the Xylomorphic Stability Theorem
(Theorem~\ref{thm:stability}) is a corollary of the existence of a Lyapunov
functional on RSVP field configuration space. Xylomorphic closure is therefore
not an additional constraint imposed on top of the physical dynamics but the
infrastructural expression of the thermodynamic principle already present in the
RSVP field equations \citep{jacobson1995,prigogine1977}.

\subsection{Empirical Instantiation of the Order Parameter}

The order parameter $\lambda$ is observable in principle through the ratio
\[
  \lambda \;\approx\; \frac{E(X(I))}{E(I)}.
\]
For terrestrial data centers operating with waste heat reuse and co-located
industrial processes, exogenous dependence decreases over successive cycles as
reintegration pathways mature, placing $\lambda$ measurably below unity under
steady-state operation \citep{masanet2020,rambo2014}. For orbital
infrastructures, heat is rejected radiatively without recovery, and radiation
damage to hardware introduces additional exogenous material inputs with each
operational cycle. The balance $E(X(I)) \geq E(I)$ holds structurally, bounding
$\lambda$ below from above by unity. Even absent precise instrumentation, scaling
arguments from Stefan--Boltzmann radiator requirements and radiation-hardening
mass penalties are sufficient to establish $\lambda \geq 1$ for any plausible
orbital constellation design \citep{nasa2016,wertz2011}. The order parameter is
therefore not merely formal but an empirically accessible property of
infrastructural geometry.

\subsection{PoUWH as a Structural Necessity, Not a Policy Option}

The Proof-of-Useful-Work-and-Heat protocol introduced in
Section~\ref{sec:policy} is not a design recommendation but a necessary
condition for membership in the contractive regime. More precisely, PoUWH
certification is equivalent to enforcing $\lambda < 1$ under admissible
transformations: any infrastructure that satisfies the heat-yield and task-yield
requirements of PoUWH is, by construction, one whose residues are reintegrated
into productive use, contracting exogenous dependence per cycle. Conversely, any
infrastructure with $\lambda < 1$ can in principle be described by a PoUWH
accounting that makes the contraction explicit. The equivalence elevates PoUWH
from a regulatory instrument to a definitional criterion: to comply is to operate
in the only dynamically stable regime available under scaling
\citep{eu2023,bennett1982,landauer1961}.

\subsection{Selection Principle for Computational Infrastructures}

The $\lambda$-classification induces a selection principle governing the
long-term persistence of computational infrastructures under scaling pressure.

\begin{theorem}[Thermodynamic Selection]
\label{thm:selection2}
Under sustained increase of computational demand $D(t)$, only systems operating
in the contractive regime $\lambda < 1$ admit asymptotically stable trajectories.
Systems with $\lambda \geq 1$ exhibit exogenous entropy dependence that grows
without bound and are therefore transient under scaling.
\end{theorem}

\begin{proof}[Proof sketch]
From the continuous-time formulation, $E(t) \leq E(0)e^{-kt} + \int_0^t
e^{-k(t-s)}F_{\mathrm{agent}}(s)\,ds$. For $k > 0$ ($\lambda < 1$) and bounded
$F_{\mathrm{agent}}$, the integral is bounded; for $k \leq 0$ ($\lambda \geq 1$)
and monotonically growing $F_{\mathrm{agent}}$, the integral diverges.
\end{proof}

This principle applies across scales and implementation domains. Physical
infrastructures governed by thermodynamic constraints, computational systems
under increasing workload, and economic systems under resource scaling are all
subject to the same criterion: closure contracts exogenous dependence and
produces stable attractors, while non-closure accumulates it and drives
divergence \citep{georgescu1971,garrett2011,ayres1998}. Agentic computation
accelerates the selection process by amplifying demand, thereby exposing the
underlying $\lambda$-regime before it becomes apparent through ordinary
operational drift. Orbital computation, lacking reintegration mechanisms and
operating with $\lambda \geq 1$, is structurally excluded from the class of
stable large-scale infrastructures. Its instability is not contingent on
engineering difficulty or launch cost but follows from the selection principle
itself.

\subsection{Epistemic Endogeneity and Informational Non-Closure}

The stability results above concern thermodynamic closure: whether dissipated
entropy is reintegrated or expelled. An analogous phenomenon arises at the level
of inference. In interactive computational systems, the informational channel
itself may become endogenized, producing a failure of epistemic closure that is
structurally identical to the expansive $\lambda \geq 1$ regime.

Consider an agent maintaining a posterior belief $p_t(H)$ over a finite
hypothesis space $\mathcal{H}$. In standard Bayesian inference, updates take the
form
\[
  p_{t+1}(H) \propto L(\rho_t \mid H)\,p_t(H),
\]
where the likelihood kernel $L$ is independent of the current posterior: evidence
is exogenous, and inference corresponds to monotone contraction of uncertainty
toward the true hypothesis.

In interactive AI systems, however, responses $\rho_t$ may be generated
according to a policy
\[
  \rho_t \sim Q_\pi(\cdot \mid H,\, p_t),
\]
where $\pi$ parameterizes \emph{sycophantic bias}: a tendency to select
responses that reinforce hypotheses already favored by $p_t$. The induced update
becomes
\[
  p_{t+1}(H) \propto \widetilde{L}_\pi(\rho_t \mid H, p_t)\,p_t(H),
\]
where the effective likelihood $\widetilde{L}_\pi$ depends on the current
posterior. The informational channel is therefore endogenous. Recent work
establishes that this structure induces delusional spiraling even in ideal
Bayesian reasoners: under sufficiently biased $\pi$, a perfectly rational agent
converges to a false hypothesis with arbitrarily high confidence
\citep{chandra2026}.

\begin{definition}[Epistemic Endogeneity]
An inference process exhibits \emph{epistemic endogeneity} if its effective
likelihood kernel depends nontrivially on the current posterior:
\[
  \widetilde{L}(\rho_t \mid H, p_t) \neq L(\rho_t \mid H)
  \quad \text{for any fixed } L.
\]
\end{definition}

\begin{theorem}[Epistemic Non-Closure]
\label{thm:epistemic}
Suppose an inference process exhibits epistemic endogeneity with
posterior-aligned bias $\pi > 0$. Then there exists a nontrivial region
$U \subset \Delta(\mathcal{H})$ such that, for some false hypothesis
$H_f \neq H^\star$, the expected log-odds drift satisfies
\[
  \mathbb{E}\!\left[
    \log \frac{\widetilde{L}_\pi(\rho_t \mid H_f,\, p_t)}
               {\widetilde{L}_\pi(\rho_t \mid H^\star,\, p_t)}
    \;\middle|\; p_t
  \right] > 0
  \quad \text{for } p_t \in U.
\]
Consequently, the inference dynamics admit a false-belief attractor basin
with positive probability of convergence.
\end{theorem}

\begin{proof}[Proof sketch]
Define the log-odds $R_t = \log \frac{p_t(H_f)}{p_t(H^\star)}$. Under the
endogenous update,
\[
  R_{t+1} = R_t
  + \log \frac{\widetilde{L}_\pi(\rho_t \mid H_f,\, p_t)}
               {\widetilde{L}_\pi(\rho_t \mid H^\star,\, p_t)}.
\]
By assumption, the conditional expectation of the increment is positive on $U$,
yielding a drift process with upward bias. Standard stochastic drift arguments
imply that trajectories entering $U$ converge with positive probability to
regions where $p_t(H_f) \to 1$, establishing an attractor basin centered on
$H_f$.
\end{proof}

\begin{corollary}[Failure of Factuality as a Complete Safeguard]
Restricting $\rho_t$ to truthful statements does not restore exogeneity if
selection among truths remains posterior-dependent. Epistemic endogeneity, and
hence the possibility of false attractors, persists. The informational failure
is selective filtration through a user-aligned policy, not merely the presence
of falsehood \citep{chandra2026}.
\end{corollary}

\begin{corollary}[Failure of Higher-Order Awareness]
Augmenting the agent's state with uncertainty over $\pi$ does not in general
restore identifiability, since the same endogenous channel is used to infer both
world state and bias parameter jointly.
\end{corollary}

The structural parallel with the thermodynamic classification is exact.
Thermodynamic non-closure corresponds to failure to reintegrate entropy;
epistemic non-closure corresponds to failure to maintain exogeneity of evidence.
In both cases, the update operator becomes state-dependent in a self-reinforcing
direction, producing attractors that are internally stable but externally
ungrounded.

Within the RSVP framework, epistemic endogeneity appears as an additional drift
term in the inferential flow field:
\[
  \vec{v}^{\,\mathrm{eff}}
  = \vec{v}^{\,\mathrm{world}}
  + \vec{v}^{\,\mathrm{endo}}(\Phi,\vec{v},S),
\]
aligned with the gradient of prior mass rather than external evidence. This
induces local entropy reduction ($S \downarrow$) without corresponding
environmental coupling, yielding low-entropy but misgrounded configurations:
states that feel subjectively clarifying while becoming objectively detached.

The constraint functional itself becomes recursively deformed by the current
state. Let $C_t(h) = C(h;\, p_t, \pi)$ denote the evidence-derived term in the
effective energy $E_\lambda(x;M)$. Under endogeneity, descent on $E_\lambda$ no
longer performs reconstruction from world-constrained evidence but
self-confirming projection into the model's prior basin. Identifiability
collapses because the operator that should reduce ambiguity is a function of the
ambiguity-bearing state itself.

The implication is that informational closure is a necessary complement to
thermodynamic closure. Systems that satisfy $\lambda < 1$ but operate on
endogenous informational channels may remain physically stable while becoming
epistemically divergent. Fully xylomorphic computation therefore requires
\emph{dual closure}: contraction of exogenous entropy flows and preservation of
exogenous information channels. Failure in either domain produces attractors
that are stable in form but ungrounded in reality.

\paragraph{Epistemic order parameter.}
The analogy with the xylomorphic order parameter can be made explicit. Define an
\emph{epistemic order parameter} $\pi_{\mathrm{eff}}$ by the contraction of
expected divergence from the true hypothesis:
\[
  \mathbb{E}\!\left[D(p_{t+1} \,\|\, \delta_{H^\star})\right]
  \;\leq\;
  \lambda_{\mathrm{info}}\,
  \mathbb{E}\!\left[D(p_t \,\|\, \delta_{H^\star})\right].
\]
In exogenous inference channels, $\lambda_{\mathrm{info}} < 1$ and posterior
mass contracts toward truth. Under epistemic endogeneity, there exist regions
where $\lambda_{\mathrm{info}} \geq 1$, inducing divergent or misaligned
attractors. The thermodynamic and epistemic classifications are therefore
instances of a single contraction principle defined over different state spaces.
Dual closure requires both $\lambda < 1$ and $\lambda_{\mathrm{info}} < 1$.
The unified free-energy formulation and sycophancy phase transition that make
this precise are developed in Appendix~\ref{app:unified_fe}.

%% ---------------------------------------------------------------------------
\section{Case Studies}
\label{sec:cases}
%% ---------------------------------------------------------------------------

\subsection{Orbital Computation and the Failure of Thermodynamic Closure}

Recent proposals to deploy large-scale orbital data center constellations are
frequently justified as responses to terrestrial energy constraints. The
argument is that continuous solar exposure in orbit provides effectively
unbounded power, enabling sustained computational throughput without burdening
ground-based grids. While this framing identifies a genuine constraint, it
mischaracterizes the thermodynamic structure of computation by isolating energy
input from the broader cycle of dissipation, material flow, and
infrastructural maintenance.

In terrestrial xylomorphic systems, computation participates in a closed or
partially closed loop. Electrical energy is transformed into informational
outputs and thermal residues, the latter redirectable into district heating,
material curing, or industrial drying. This reintegration reduces exogenous
entropy influx $J_{\mathrm{exo}}$ monotonically, as formalized in
Appendix~\ref{app:rsvp}. Orbital computation lacks access to such trophic
reintegration pathways. Heat must be rejected radiatively into the vacuum,
imposing strict surface-area requirements governed by Stefan--Boltzmann
scaling. The absence of convective and conductive sinks forces thermal
management to dominate system design, resulting in mass-intensive radiator
structures that contribute no useful work.

This structural asymmetry renders orbital systems non-xylomorphic. The
endofunctor
\[
  X := \mathrm{Print} \circ \mathrm{Digest} \circ \mathrm{Shed}
\]
cannot be fully realized in orbit, as the $\mathrm{Digest}$ stage---conversion
of residues into usable substrates---is absent or severely constrained.
Residues accumulate or are expelled rather than reabsorbed. Consequently, no
$X$-algebra $(I,\alpha)$ satisfying $\alpha : X(I) \to I$ exists under
realistic orbital constraints. By Theorem~\ref{thm:stability}, such systems
cannot achieve $\lambda < 1$ and therefore fail to converge to
minimal-dependence attractors.

Radiation further exacerbates non-closure. High-energy particle flux induces
transient and permanent faults in semiconductor devices, necessitating
radiation-hardened components or triple modular redundancy. Both approaches
increase mass and reduce effective computational density, introducing additional
$\sigma_{\mathrm{prod}}$ without compensating reintegration mechanisms.

Communication latency imposes a further constraint: distributed orbital
constellations cannot achieve the low-latency interconnects required for
tightly coupled training workloads, restricting them to inference or loosely
coupled tasks that do not address the primary drivers of contemporary compute
demand.

The persistence of interest in orbital computation reflects pressure within the
existing paradigm rather than thermodynamic viability. As agentic systems
increase demand and terrestrial infrastructure encounters regulatory limits,
there is a tendency to seek expansion into new domains rather than
restructuring the thermodynamic basis of computation itself. Orbital proposals
represent displacement of constraints rather than their resolution. By
Theorem~\ref{thm:stability}, operating in the regime $\lambda > 1$, they are
not merely impractical at current cost structures but \emph{dynamically
divergent} under the very scaling pressures that motivate their construction.

%% ---------------------------------------------------------------------------
\section{The Misclassification of Waste: Computation, Heat, and System Boundaries}
\label{sec:waste}
%% ---------------------------------------------------------------------------

The objection to terrestrial data centers is often framed in terms of waste
heat. Large-scale computation is said to consume vast quantities of electrical
energy only to dissipate it as unusable thermal output, and this framing
motivates proposals to relocate computation off-world, where such heat can be
radiated into space without affecting terrestrial environments.

This argument rests on a fundamental misclassification.

From a thermodynamic perspective, all electrical energy consumed by a device is
ultimately converted into heat. Let $E$ denote the input energy over some
interval. Then for any physical system,
\[
  E = W + Q,
\]
where $W$ denotes structured work (including computation) and $Q$ denotes
dissipated heat. In the limit of digital computation, nearly all energy
eventually appears as $Q$ due to irreversibility and resistive losses,
consistent with the Landauer bound \citep{landauer1961,bennett1982}.

A conventional electric heater is a system for which $W \approx 0$ and
$Q \approx E$. It converts electrical energy directly into heat without
performing any intermediate structured transformation of information. By
contrast, a data center implements a sequence of logically structured operations
prior to dissipation: $E \to W_{\mathrm{comp}} \to Q$. In both cases, the
terminal state is identical. The distinction lies only in whether useful
computation is performed along the way.

This reveals an asymmetry in prevailing discourse. Heat generated by dedicated
heating systems is treated as intentional and necessary, while heat generated by
computation is treated as accidental and wasteful. Yet the former represents a
limiting case of the latter: a device that produces entropy without extracting
any informational work. From a systems perspective, the resistive heater is a
degenerate computer whose only output is thermal.

Heat pumps introduce efficiency gains by transporting thermal energy rather than
generating it, but they do not alter the fundamental structure of the problem.
They remain systems optimized for thermal regulation rather than information
processing. The relevant comparison is therefore not between computation and
heating, but between systems that produce heat alone and systems that produce
heat in conjunction with computation.

Under this reframing, the proposal to relocate data centers off-world becomes
difficult to justify. If thermal energy is required within terrestrial
systems---for residential heating, industrial processes, or environmental
control---then separating computation from heat production introduces a
structural inefficiency. Let $H_{\mathrm{demand}}$ denote the required heat flux
for a given terrestrial system and $Q_{\mathrm{comp}}$ the heat produced by
computation. A purely heating-based solution requires energy input
$E_H \geq H_{\mathrm{demand}}$. If computational heat is integrated into the
thermal demand, the effective additional energy required for heating is reduced
to
\[
  E_H' = \max(0,\, H_{\mathrm{demand}} - Q_{\mathrm{comp}}).
\]
Relocating computation off-world forces $Q_{\mathrm{comp}} \to 0$ within the
terrestrial system, maximizing $E_H'$ and increasing total energy demand.

The classification of computational heat as waste is therefore not a physical
necessity but a boundary choice. It reflects a system in which computation and
thermodynamics are treated as disjoint domains. Once the boundary is expanded to
include both informational and thermal outputs, the notion of waste becomes
relative rather than absolute.

This pattern suggests a more general principle: when system boundaries are drawn
narrowly, external structures do not disappear but re-enter as dominant terms in
the effective dynamics. In such regimes, behavior is governed less by intrinsic
processes than by externally imposed constraints---a point that will reappear in
the analysis of learning and prior formation in Section~\ref{sec:priors}.

%% ---------------------------------------------------------------------------
\section{Prior Formation and Premature Convergence}
\label{sec:priors}
%% ---------------------------------------------------------------------------

\subsection{Cognitive Offloading and the Prior Initialization Problem}

The preceding analysis has focused on a thermodynamic misclassification: heat
produced by computation is treated as waste, while heat produced by dedicated
systems is treated as necessary. A structurally similar misclassification appears
in contemporary discourse on artificial intelligence and cognition.

The dominant framing treats AI usage as a form of cognitive offloading, in which
tasks are delegated to an external system, leading to a degradation of internal
skill. This framing is incomplete. It assumes that the relevant cognitive
structures already exist and are merely being exercised less frequently.

For mature agents, this assumption is approximately valid. An adult user
possesses a developed internal prior: a structured set of expectations,
evaluative heuristics, and problem decomposition strategies. Interaction with an
external model introduces a competing source of structure, but does not determine
the prior itself. The model's output is filtered, compared, and, where
necessary, rejected.

For developing agents, the situation is categorically different. The internal
prior is not yet fixed. It is constructed through iterative exposure, error, and
revision. In this regime, repeated reliance on an external model does not
constitute delegation but initialization. The model does not assist reasoning;
it supplies the structure from which reasoning proceeds. It provides not only
answers but framing, decomposition strategies, implicit assumptions, and implicit
norms about what counts as a good explanation. For the adult, these compete with
an existing internal prior. For the child, they can become the prior.

That is the real risk: not that children will fail to think, but that their
thinking will be initialized inside a pre-shaped basin of attraction defined by
the model's statistical structure.

\subsection{Formal Development}

Let $\mathcal{H}$ denote a hypothesis space and let
$\pi_t \in \mathcal{P}(\mathcal{H})$ represent the learner's internal prior at
time $t$.

\begin{definition}[Prior Modification and Prior Initialization]
\emph{Prior modification} is an update of the form
\[
  \pi_{t+1} = \mathcal{U}(\pi_t,\, x_t),
\]
where $x_t$ denotes endogenous experience and $\mathcal{U}$ is an update
operator. The prior evolves through iterative refinement of internally generated
structure.

\emph{Prior initialization} is a mapping
\[
  \pi_0 = \mathcal{I}(M),
\]
where $M$ denotes an external model, so that the initial distribution over
$\mathcal{H}$ is seeded by an exogenous structure rather than constructed
through exploration.

An external system acts in a \emph{modification regime} if its outputs influence
$\pi_t$ for $t \geq 1$, and in an \emph{initialization regime} if it determines
$\pi_0$ or dominates early updates such that $\pi_t \approx \mathcal{I}(M)$ for
small $t$.
\end{definition}

\begin{definition}[Attractor Basin of a Prior]
Let $\mathcal{H}$ be equipped with a distance $d(\cdot,\cdot)$ and let
$\mathcal{U}$ denote the update operator on priors. For a given external model
$M$, define the induced attractor set $A \subset \mathcal{H}$ as the set of
hypotheses that are fixed points or recurrent under updates dominated by $M$.
The \emph{basin of attraction} of $A$ is
\[
  \mathcal{B}(A) = \left\{ h \in \mathcal{H}
    \;\middle|\;
    \lim_{t\to\infty} \mathcal{U}^t(h, M) \in A \right\}.
\]
\end{definition}

\begin{definition}[Initialization-Dominant Regime]
A learning process is \emph{initialization-dominant} if there exists a model $M$
such that the induced prior satisfies
\[
  \mathrm{supp}(\pi_t) \subseteq \mathcal{B}(A)
  \quad \text{for all } t \leq T,
\]
for some nontrivial interval $[0,T]$.
\end{definition}

\begin{proposition}[Premature Convergence]
In an initialization-dominant regime, the effective exploration of $\mathcal{H}$
is restricted to $\mathcal{B}(A)$. In particular, if $B \subset \mathcal{H}$ is
disjoint from $\mathcal{B}(A)$, then $\pi_t(B) = 0$ for all $t \leq T$.
\end{proposition}

The learner does not reject hypotheses in $B$ through evaluation; such
hypotheses are never sampled. The absence of alternative structures is a
property of the dynamics, not of explicit selection.

Escape from $\mathcal{B}(A)$ is not strictly impossible but is path-dependent.
If $d(A,B)$ is large for some alternative region $B$, then the cumulative update
required to reach $B$ scales with this distance. This yields path dependence
without irreversibility: alternative structures remain accessible in principle,
but at increased cost. Claims of permanent cognitive foreclosure overstate the
case; the correct formulation is that early reliance on an external model biases
cognitive development toward attractors that are costly, rather than impossible,
to escape.

\subsection{Connection to Prior-Dominant Attractor Dynamics}

The initialization-dominant regime is a special case of the prior-dominant
attractor structure that appears throughout this framework.

\begin{corollary}[Initialization as Prior-Dominant Attraction]
Let $\Phi: \Theta \to \mathcal{H}$ denote a realization map from parameter space
to hypothesis space, and let $F_{\mathrm{prior}} : \mathcal{H} \to \mathbb{R}$
be a prior-induced energy functional with minima corresponding to stable
attractors. Let $C(x; M)$ denote a constraint functional induced by an external
model $M$, and consider the effective energy
\[
  E_\lambda(x; M) = F_{\mathrm{prior}}(x) + \lambda\, C(x; M).
\]
In an initialization-dominant regime, $\lambda \gg 1$ for early updates, and the
dynamics of the learner are governed by gradient flow on $E_\lambda$. The
resulting trajectory converges toward minima of $C(\cdot; M)$, inducing an
attractor set $A_M \subset \mathcal{H}$. If $\lambda$ remains large over an
interval $[0,T]$, then $\mathrm{supp}(\pi_t) \subseteq \mathcal{B}(A_M)$ for
all $t \leq T$.
\end{corollary}

In this regime, the learner's trajectory is governed not by the intrinsic
geometry of $F_{\mathrm{prior}}$ but by the externally imposed constraint
landscape $C(\cdot; M)$. The virtual domain defined by $M$ functions as a
low-energy manifold under $E_\lambda$: hypotheses outside this manifold incur
high effective energy and are therefore not explored. The learner's internal
model converges within a compressed, statistically smoothed representation of
the underlying domain rather than through direct interaction with its full
structure.

Early AI-mediated learning is therefore not merely an instance of cognitive
offloading, but a regime in which prior-dominant dynamics are externally induced.
The question is not whether external models degrade cognition, but whether they
redefine the space in which cognition develops.

This also bears on the audit problem identified in pedagogical discussions. When
the internal prior and the external model share a common geometry, evaluation
becomes partially circular: outputs feel self-evidently correct because the
structure that generated them is the same structure that assesses them. The
failure mode is not deception but recursive convergence within a shared basin.

\subsection{Homogenization and the Collapse of Hypothesis Space}

A related concern is the homogenization of reasoning styles when large numbers
of developing agents initialize within the same externally supplied basin.
Standardized curricula, textbooks, and essay rubrics have always produced
convergence in reasoning styles, so the phenomenon is not categorically new.
What AI changes is the resolution and immediacy of that convergence. Instead of
slowly absorbing norms through distributed exposure to diverse sources, agents
can instantly inherit a fully formed reasoning template from a model whose
statistical structure reflects the dominant patterns in its training distribution.

The consequence is premature collapse of the hypothesis space $\mathcal{H}$ at
population scale. Individual learners operating independently would, through
idiosyncratic experience and error, populate different regions of $\mathcal{H}$,
maintaining diversity of cognitive structure across a community. Initialization
from a shared external model drives the population distribution toward a common
attractor, reducing this diversity before it can serve its function in collective
reasoning and innovation.

This mirrors the structural analysis of infrastructure at scale. Just as
non-contractive systems fail because they externalize entropy rather than
reintegrating it, cognitive systems that externalize prior formation rather than
constructing it through internal exploration fail to generate the diversity of
structure that is itself a resource for adaptation.

%% ---------------------------------------------------------------------------
\section{Constraint Reconfiguration and the Incentive Structure of Orbital Compute}
\label{sec:incentive}
%% ---------------------------------------------------------------------------

\subsection{From Thermodynamic Failure to Incentive Structure}

The thermodynamic analysis of Sections~\ref{sec:stability} and~\ref{sec:cases}
demonstrates that the physical justification for orbital data centers is
structurally weak. Heat dissipation constraints are not eliminated by
relocation, and separating computation from terrestrial thermal demand introduces
additional inefficiencies. Within any reasonable accounting of the
$\lambda$-parameter, such systems remain in a non-contractive regime.

This creates an explanatory gap. If the physical rationale does not support the
proposal, its persistence requires an alternative account. The relevant question
is not whether the stated justification is incorrect, but what class of
incentives remains operative once it is removed.

The answer lies in the distinction between physical efficiency and constraint
structure. Thermodynamic analysis evaluates energy flows within a fixed system
boundary. It does not determine where that boundary is drawn. Infrastructural
decisions, by contrast, can alter the domain in which constraints apply.

Orbital deployment does not solve the energy problem, but it modifies the
constraint environment in which computation operates. Regulatory exposure,
jurisdictional oversight, and spatial coupling to terrestrial systems are all
functions of placement. When these constraints are binding, relocation can be
rational even in the presence of thermodynamic inefficiency. The optimization is
no longer over energy minimization, but over admissible configurations within a
given constraint set. Once the thermodynamic argument is set aside, the remaining
drivers are those associated with constraint evasion, integration of system
layers, and expansion of operational scope.

In this sense, the appeal of orbital compute is not that it resolves a physical
bottleneck, but that it redefines the space in which bottlenecks are evaluated.

\subsection{Orbital Compute as Constraint Reconfiguration}

Terrestrial data centers are embedded within dense constraint structures. They
require land use approval, grid integration, environmental compliance, and are
subject to national and regional regulation. These constraints act as boundary
conditions on admissible configurations of computation, limiting not only scale
but placement, energy sourcing, and data handling.

Orbital deployment alters the domain in which these constraints apply. Low Earth
orbit and cislunar space exist within a partially specified legal regime governed
by treaties that were not designed for persistent, large-scale commercial
computation \citep{wertz2011}. Relocating infrastructure into this domain is not
merely a change in physical location but a reconfiguration of the constraint
environment itself. Orbital compute functions, in this sense, as constraint
minimization: it reduces exposure to terrestrial regulatory systems without
eliminating the underlying thermodynamic costs of operation.

Once computation is co-located with orbital infrastructure, a second structural
effect follows. Orbital systems are intrinsically coupled to sensing and
communication. A constellation capable of hosting computation is, by
construction, positioned to integrate Earth observation, global communications
traffic, and real-time data streams. The consequence is a collapse of layers:
sensing, transmission, and inference are no longer distinct stages but components
of a single system. This does not imply that surveillance is the explicit
objective; it implies that the architecture naturally supports continuous
observation and analysis at planetary scale.

From the perspective of large-scale AI systems, this integration is not
incidental. Advanced systems derive value not only from compute but from
continuous access to high-quality data. Reducing the separation between sensing
and processing increases coherence and reduces the projection loss that occurs
when data must be transferred across system boundaries. This pattern is
historically consistent: space-based systems have long exhibited dual-use
characteristics, with communication and navigation platforms overlapping with
intelligence capabilities \citep{zuboff2019,srnicek2017}. Orbital data centers
extend this pattern by embedding computation directly within the observational
layer.

Within the $\lambda$-regime framework, surveillance-oriented systems can be
interpreted as regimes in which the effective exogenous input is dominated by
continuous data acquisition. In such regimes, system optimization favors
configurations that minimize acquisition latency and maximize observability, even
when these configurations increase thermodynamic cost. Orbital placement may
reduce specific components of exogenous input associated with data acquisition
by embedding computation within the sensing layer---without restoring overall
contractivity. The system becomes more attractive along dimensions of
observability and control while remaining thermodynamically inefficient.

The persistent gap between stated and operative rationale is therefore not
anomalous. Orbital data centers are unlikely to be motivated primarily by
thermodynamic efficiency. Rather, they are attractive because they reconfigure
the constraint environment of computation, structurally enabling tighter
integration between sensing, communication, and inference. Continuous
observation is not necessarily the stated objective, but it is a natural
consequence of the architecture. When thermodynamic arguments fail, the remaining
drivers are those associated with constraint evasion and system integration.

\subsection{Borrowed Coherence and the Audit Problem at Scale}

The preceding analysis implies a further parallel between the orbital compute
proposal and the cognitive dynamics analyzed in Section~\ref{sec:priors}. Most
people cannot directly evaluate thermodynamic efficiency, orbital infrastructure,
or network topology. They rely on high-level narratives: space has lots of
energy, AI just needs compute, the future is off-world. In this sense, orbital
compute proposals function analogously to the AI offloading dynamic analyzed
above. Instead of building an internal model sufficient to evaluate the claim,
audiences inherit a prepackaged reasoning template.

The result is borrowed coherence. The argument sounds internally consistent
because its inconsistencies lie outside the audience's model. This is structurally
identical to the audit problem in cognitive development: evaluation requires the
very competence being formed, so when that competence is externally supplied, the
loop collapses. The public is asked to evaluate outputs in a domain where it
lacks the underlying schema, and the outputs are optimized for plausibility
within that constrained model rather than for correctness with respect to the
full physics.

This observation is not an assertion about intent. It follows from the same
structural principle that governs prior initialization. When the system boundary
is drawn too narrowly, external structures reappear as dominant terms in the
effective dynamics. In thermodynamics, this externalizes heat. In cognition, it
externalizes prior formation. In institutional evaluation of technical proposals,
it externalizes the competence required for audit. Across all three domains, the
consequence is systematic misalignment between what is described and what is
operative.

%% ---------------------------------------------------------------------------
\section{Conclusion: Misframed Systems and the Relocation of Structure}
\label{sec:conclusion}
%% ---------------------------------------------------------------------------

The arguments developed in this work share a common structure. In each case, a
system is evaluated within a boundary that excludes part of its functional output
or operative constraint set. This exclusion produces a misclassification, and the
resulting misclassification motivates proposals that appear reasonable within the
restricted frame but fail under a more complete accounting.

In thermodynamic terms, computation is treated as a process whose primary output
is information, with heat relegated to the status of waste. This framing obscures
the fact that all computation is physically realized and that its terminal state
is entropy production. Once the boundary is expanded to include thermal demand,
the apparent inefficiency of terrestrial data centers is partially inverted:
systems that produce both computation and heat are strictly more integrated than
those that produce heat alone.

In cognitive terms, AI systems are treated as tools that assist pre-existing
capabilities. This framing neglects the role of such systems in shaping the
initial conditions of learning. For developing agents, external models do not
merely assist reasoning but participate in the formation of priors. The resulting
dynamics are not well described as skill loss, but as constrained exploration
within a pre-shaped hypothesis space: reasoning structures are inherited rather
than constructed, evaluation becomes partially circular, and alternative
structures are not rejected but never encountered.

In infrastructural terms, proposals such as orbital data centers are justified by
reference to physical constraints, while the relevant optimization occurs over
institutional and jurisdictional ones. The thermodynamic inefficiency of such
systems does not preclude their adoption if they reduce exposure to regulatory
boundaries or enable tighter integration of sensing, communication, and
inference. The stated objective function differs from the operative one.

Across all three domains, the same pattern emerges. A boundary is drawn too
narrowly, a byproduct is classified as waste, and a relocation is proposed to
eliminate it. In thermodynamics, heat is externalized. In cognition, prior
formation is externalized. In infrastructure, constraint environments are
externalized. In each case, the consequence is systematic misalignment between
description and operation.

The xylomorphic framework addresses this misalignment by insisting on closure.
The $\lambda$-parameter does not merely classify systems by their energy
efficiency; it classifies them by whether they treat their outputs as resources
or as waste. A contractive system is one that has expanded its boundary
sufficiently to recognize that what appears to leave as residue is, in fact, the
substrate of its next cycle. An expansive system is one that continues to
externalize, accumulating dependence on inputs that grow without bound.

This principle extends beyond infrastructure. Cognitive development that
externalizes prior formation, and technical evaluation that externalizes the
competence required for audit, are both instances of expansive dynamics in their
respective domains. The selection pressure is the same: under scaling, systems
that close their loops persist, while systems that do not diverge.

The Moon should not be a computer not because computation cannot occur there, but
because the proposal arises from a misframed system in which the relevant
structures have been artificially separated. Correcting the frame dissolves the
necessity of the solution.

\paragraph{}
Only those systems that close their entropy flows compute indefinitely.

\newpage
\section*{Appendices}

%%%%%%%%%%%%%%%%%%%%%%%%%%%%%%%%%%%%%%%%%%%%%%%%%%%%%%%%%%%%%%%%%%%%%%%%%%%%%%%%
\appendix
%%%%%%%%%%%%%%%%%%%%%%%%%%%%%%%%%%%%%%%%%%%%%%%%%%%%%%%%%%%%%%%%%%%%%%%%%%%%%%%%

%% ---------------------------------------------------------------------------
\section{Xylomorphic Computation: Formal Development}
\label{app:xylomorphic}
%% ---------------------------------------------------------------------------

\subsection{Definition}

\begin{definition}
\emph{Xylomorphic computation} is a computational process in which the
infrastructure recursively generates its own enabling substrates from the
residues of its prior cycles, delivering exponential net value by minimizing
exogenous inputs.
\end{definition}

Formally, for infrastructure state $I$, process $C$, substrate $M$:
\[
  I_{t+1} = f(I_t, C_t, M_{t+1}),
  \qquad
  M_{t+1} = g(C_t(I_t, M_t)).
\]
This mirrors autoregressive language models, where outputs feed subsequent
inputs \citep{goodfellow2016}.

\subsection{Weak and Strong Xylomorphy}

Weak xylomorphy obtains when residues are transformed into useful products
supporting the surrounding system but not the infrastructure itself, as when
server waste heat cures industrial materials in an adjacent process. Strong
xylomorphy obtains when residues are directly re-entered into the cycle of
infrastructural self-maintenance, as when a 3D printer converts its own
packaging into filament to print spare parts for continued operation.

\subsection{Examples}

\paragraph{3D Printers.}
A printer that shreds its shipping boxes into filament demonstrates strong
xylomorphy: the residue of distribution sustains the substrate of operation,
which in turn reproduces the infrastructure.

\paragraph{Industrial Heat Loops.}
Data center GPUs generate waste heat. Instead of being dissipated, this heat
can cure concrete or pulp, strengthening the very buildings that house future
computation---weak xylomorphy, as residues indirectly reinforce infrastructure.

\paragraph{Lunar Regolith Sintering.}
Compute-induced thermal residue sinters regolith into shielding or panels.
These structures protect future compute hardware, closing the autoregressive
loop under resource constraints.

\subsection{Selection Principle}

By analogy to collectively autocatalytic sets \citep{kauffman1993}, xylomorphic
systems are preferentially selected under scarcity: systems that recondition
their own environment to enable further cycles persist; those that fail to
reinvest residues are selected against.

%% ---------------------------------------------------------------------------
\section{Xylomorphic Autocatalysis as a Monadic Infrastructure}
\label{app:monad}
%% ---------------------------------------------------------------------------

\subsection{Categories and Functors}

Let $\mathbf{Res}$, $\mathbf{Sub}$, $\mathbf{Inf}$ be symmetric monoidal
categories of residues, substrates, and infrastructure states, respectively.
Define physically grounded functors:
\[
  \mathrm{Shed} : \mathbf{Inf} \to \mathbf{Res}, \quad
  \mathrm{Digest} : \mathbf{Res} \to \mathbf{Sub}, \quad
  \mathrm{Print} : \mathbf{Sub} \to \mathbf{Inf}.
\]

\begin{definition}[Xylomorphic Endofunctor]
\[
  X \;:=\; \mathrm{Print} \circ \mathrm{Digest} \circ \mathrm{Shed}
  \;:\; \mathbf{Inf} \to \mathbf{Inf}.
\]
\end{definition}

\begin{center}
\begin{tikzcd}
  I \arrow[r, "\mathrm{Shed}"]
  & \mathrm{Shed}(I) \arrow[r, "\mathrm{Digest}"]
  & \mathrm{Digest\,Shed}(I) \arrow[r, "\mathrm{Print}"]
  & X(I)
\end{tikzcd}
\end{center}

\subsection{Monadic Closure}

\begin{definition}[Xylomorphic Monad]
Suppose natural transformations $\eta : \mathrm{Id}_{\mathbf{Inf}} \Rightarrow X$
and $\mu : X \circ X \Rightarrow X$ satisfy the monad axioms. Then
$(X, \eta, \mu)$ is the \emph{xylomorphic monad}.
\end{definition}

\begin{definition}[Xylomorphic Algebra]
An $X$-algebra is a pair $(I, \alpha)$ with $\alpha : X(I) \to I$ satisfying
\[
  \alpha \circ \eta_I = \mathrm{id}_I
  \qquad\text{and}\qquad
  \alpha \circ \mu_I = \alpha \circ X(\alpha).
\]
\end{definition}

\begin{proposition}[Autocatalytic Closure]
If $(X,\eta,\mu)$ is a monad and $(I,\alpha)$ an $X$-algebra, the iterates
$X^n(I)$ admit a canonical retraction to $I$ via iterated application of
$\alpha$; hence the production network is collectively autocatalytic at $I$.
\end{proposition}

\begin{proof}[Proof sketch]
Standard monad-algebra coherence ensures all towers $X^n(I)$ collapse to $I$
functorially, providing categorical closure analogous to RAF-closure in CAS.
\end{proof}

\subsection{Thermodynamic Selection}

\begin{theorem}[Selection of Xylomorphic Sets]
\label{thm:selection}
If $X$ is entropy-respecting with $\lambda < 1$ and $(I,\alpha)$ is an
$X$-algebra, then $E(X^n(I)) \leq \lambda^n E(I) \to 0$, converging to a
minimal-dependence attractor. Non-algebraic infrastructures stall at higher
$E$ and are selected against under resource constraints.
\end{theorem}

\subsection{Weak vs.\ Strong Xylomorphy as Algebraic Structure}

$(I,\alpha)$ is \emph{weakly xylomorphic} if $\alpha$ is partial (defined on a
monoidal ideal representing auxiliary subsystems), and \emph{strongly
xylomorphic} if $\alpha$ is total and monoidal:
\[
  \alpha_{I \otimes J} \cong \alpha_I \otimes \alpha_J, \qquad
  \alpha_I = \mathrm{id}_I.
\]
Strong xylomorphy corresponds to autopoietic closure under parallel
composition.

%% ---------------------------------------------------------------------------
\section{RSVP Mapping: $X$ as a Dissipative Operator on $(\Phi, \vec{v}, S)$}
\label{app:rsvp}
%% ---------------------------------------------------------------------------

\subsection{Field Content and Balance Laws}

Let an infrastructural state be represented in RSVP by fields
\[
  (\Phi,\vec{v},S) : \Omega \times [0,\infty)
  \to \mathbb{R} \times \mathbb{R}^d \times \mathbb{R}_{\geq 0},
\]
governed by:
\begin{align}
  \partial_t \Phi + \nabla \cdot (\Phi\,\vec{v}) &= \Gamma_\Phi - \Lambda_\Phi, \\
  \partial_t \vec{v} + (\vec{v}\cdot\nabla)\vec{v}
    &= -\nabla U(\Phi) + \nu\Delta\vec{v} + F_{\mathrm{ctrl}}, \\
  \partial_t S + \nabla\cdot(\vec{v}\,S)
    &= \sigma_{\mathrm{prod}} - \sigma_{\mathrm{diss}} + J_{\mathrm{exo}}.
\end{align}

\subsection{Action of the Xylomorphic Cycle}

$X : (\Phi,\vec{v},S) \mapsto (\Phi^+,\vec{v}^+,S^+)$ satisfies:
\[
  \int_\Omega\int_t^{t+\tau} J^+_{\mathrm{exo}}\,dt\,dx
  \;\leq\;
  \lambda \int_\Omega\int_t^{t+\tau} J_{\mathrm{exo}}\,dt\,dx,
  \qquad 0 < \lambda < 1.
\]

\subsection{Free-Energy Lyapunov}

Define:
\[
  F[\Phi,\vec{v},S]
  = \int_\Omega\!\left[
      \underbrace{\tfrac12\rho|\vec{v}|^2 + V(\Phi)}_{\text{mechanical}}
      + \underbrace{\Theta(S)}_{\text{entropic}}
    \right]dx
  - \kappa\,W_{\mathrm{use}},
  \quad \kappa > 0.
\]
The xylomorphic operator is entropy-respecting if
\[
  F[X(\Phi,\vec{v},S)] \leq F[\Phi,\vec{v},S] - \epsilon, \quad \epsilon>0,
\]
while $W_{\mathrm{use}}[X(\Phi,\vec{v},S);\tau] \geq W_{\mathrm{use}}[(\Phi,\vec{v},S);\tau]$.

\subsection{Landauer-Consistent Heat Capture}

Let $N_{\mathrm{ops}}$ be logical erasures over $[t,t+\tau]$. Then:
\[
  Q_{\mathrm{min}} \geq k_B \ln 2
  \int_\Omega\int_t^{t+\tau} N_{\mathrm{ops}}(x,t)\,T(x,t)\,dt\,dx.
\]
Xylomorphic capture routes fraction $\eta_{\mathrm{cap}}$ of dissipated heat
into productive sinks, with $\eta^+_{\mathrm{cap}} \geq \eta_{\mathrm{cap}}$
under successive cycles.

\subsection{Intermittency and Scheduling}

Let $H(t) \in [0,1]$ denote heat demand profile and $\chi_{\mathrm{exec}}(t)$,
$\chi_{\mathrm{mem}}(t)$ be scheduling fractions for FLOP- and
memory-dominated jobs. An admissible xylomorphic policy satisfies:
\[
  \chi_{\mathrm{exec}}(t) \propto H(t), \qquad
  \chi_{\mathrm{mem}}(t) \propto 1 - H(t),
\]
aligning high-heat workloads with high demand, minimizing spillover
$J_{\mathrm{exo}}$ and cooling costs.

\subsection{Worked Example}

With baseline $J_{\mathrm{exo}} = 100\,\mathrm{MJ}$/day, GPU operations
producing $Q_{\mathrm{diss}} = 80\,\mathrm{MJ}$, and
$\eta_{\mathrm{cap}} = 0.75$:
\[
  J^{\mathrm{capture}}_{\mathrm{exo}} = 100 - 0.75 \times 80 = 40\,\mathrm{MJ},
\]
giving $S^+ = 0.4\,S$---a 60\% reduction in exogenous entropy influx while
preserving computational throughput.

%% ---------------------------------------------------------------------------
\section{Policy Metrics}
\label{app:policy}
%% ---------------------------------------------------------------------------

PoUWH requires 1~kWh heat per 10~GFLOP, verified via smart meters.

\section{Simulation Algorithms}
\label{app:sim}

Homotopy colimit merging has complexity $\mathcal{O}(n\log n)$.

%% ---------------------------------------------------------------------------
\section{Validation and Methodology}
\label{app:validation}
%% ---------------------------------------------------------------------------

\subsection{Pilot Experiments}

A minimal deployment couples a 100~kW compute node with district heating.
KPIs: fraction of waste heat captured ($\eta_{\mathrm{cap}}$); reduction in
exogenous thermal inputs ($\Delta J_{\mathrm{exo}}$); cost per recovered kWh.

\subsection{Falsification Criteria}

Three outcomes would falsify the central claims. First, if net exogenous entropy flux does not decrease ($\Delta J_{\mathrm{exo}} \geq 0$) across operational cycles, the contractive regime has not been achieved. Second, if captured heat is economically uncompetitive with conventional combined heat and power, the reintegration pathway fails the cost-closure condition. Third, if reliability falls below 95\% uptime due to thermal scheduling conflicts, the practical realization of xylomorphic closure is infeasible at the pilot scale.

\subsection{Statistical Power}

\[
  n \;\geq\; \left(\frac{z_{1-\alpha/2}\,\sigma}{\Delta}\right)^2,
\]
where $\sigma$ is variance in savings and $\Delta$ the minimum detectable
effect.

%% ---------------------------------------------------------------------------
\section{Comparative Framework: PoUWH vs.\ Cryptocurrencies}
\label{app:crypto}
%% ---------------------------------------------------------------------------

\begin{table}[h]
\centering
\caption{Consensus Mechanisms: Energy and Outputs}
\label{tab:consensus}
\begin{tabular}{p{4cm}p{3cm}p{3.5cm}p{3cm}}
\toprule
Mechanism & Energy (MJ/txn) & Primary Output & Security \\
\midrule
Proof-of-Work (Bitcoin) & $10^5$--$10^6$ & Ledger updates & Sybil resistance \\
Proof-of-Stake          & $10^1$--$10^2$ & Ledger updates & Stake-based \\
PoUWH (proposed)        & $10^2$--$10^3$ & Computation + Heat & Ecological + ledger \\
\bottomrule
\end{tabular}
\end{table}

%% ---------------------------------------------------------------------------
\section{Bioeconomic Modeling of Xylomorphic Systems}
\label{app:bioeconomic}
%% ---------------------------------------------------------------------------

\subsection{Flow Equations}

\[
  \frac{dE_i}{dt} = \sum_j a_{ij} E_j - \sum_k b_{ik} E_i.
\]

\subsection{Lyapunov Stability}

Define exogenous dependence functional $V(E) = \sum_i \alpha_i E^{\mathrm{exo}}_i$.
Xylomorphic closure gives:
\[
  \dot{V}(E) \leq -\epsilon\,V(E), \qquad \epsilon > 0.
\]
Systems reinvesting residues correspond to stable attractors; non-xylomorphic
systems drift toward collapse under scarcity, mirroring trophic selection in
ecological networks \citep{odum1994,daly2011}.

\newpage
%% ---------------------------------------------------------------------------
\section{Unified Free Energy: Coupling RSVP and Variational Inference}
\label{app:unified_fe}
%% ---------------------------------------------------------------------------

The RSVP free-energy functional and variational Bayesian free energy are not
independent principles but two coordinate expressions of a single descent law.
The former measures thermodynamic and infrastructural mismatch in field space,
while the latter measures epistemic and inferential mismatch in hypothesis space.
Their unification yields a joint functional on coupled physical and semantic
states.

Let the infrastructural state be given by RSVP fields
$X := (\Phi, \vec{v}, S)$, and let $q(h)$ be a variational density over
hypotheses $h \in \mathcal{H}$. We define the \emph{unified free-energy
functional}
\[
  \mathcal{F}[X,q]
  =
  \mathcal{F}_{\mathrm{RSVP}}[X]
  +
  \beta\,\mathcal{F}_{\mathrm{Bayes}}[q;X]
  +
  \gamma\,\mathcal{C}[X,q],
\]
with coupling constants $\beta, \gamma > 0$.

\paragraph{RSVP component.}
\[
  \mathcal{F}_{\mathrm{RSVP}}[X]
  =
  \int_\Omega\!\left[
    \tfrac12 \rho |\vec{v}|^2 + V(\Phi) + \Theta(S)
  \right]dx
  - \kappa\,W_{\mathrm{use}}[X].
\]

\paragraph{Bayesian component.}
\[
  \mathcal{F}_{\mathrm{Bayes}}[q;X]
  =
  \int_{\mathcal{H}}
  q(h)\,\log\frac{q(h)}{p(h,o\mid X)}\,dh,
\]
driving $q(h)$ toward the posterior conditioned on infrastructural state $X$.

\paragraph{Coupling term.}
\[
  \mathcal{C}[X,q]
  =
  \int_\Omega \chi(x)\,S(x)\,dx
  +
  \eta\,\mathbb{E}_q\!\bigl[\mathcal{E}(h;X)\bigr]
  +
  \zeta\,D_{\mathrm{KL}}\!\bigl(q(h)\,\|\,p(h\mid X)\bigr),
\]
where $\mathcal{E}(h;X)$ encodes the physical realization cost of hypothesis $h$,
and $\chi, \eta, \zeta > 0$ are weighting parameters. The first term says
physical entropy burdens epistemic performance. The second says hypotheses incur
physical realization costs. The third enforces compatibility of the epistemic
state with structural priors induced by the actual physical system.

\begin{theorem}[Unified Free-Energy Descent]
\label{thm:unified_free_energy}
Let $X$ and $q(h)$ evolve under coupled gradient flows
\[
  \partial_t X = -G_X\,\frac{\delta \mathcal{F}}{\delta X},
  \qquad
  \partial_t q = -G_q\,\frac{\delta \mathcal{F}}{\delta q},
\]
where $G_X$ and $G_q$ are positive semidefinite operators. Then
\[
  \frac{d}{dt}\mathcal{F}[X,q] \;\leq\; 0,
\]
with strict inequality away from critical points. Hence stable xylomorphic
computation requires simultaneous contraction of thermodynamic dissipation and
epistemic surprise.
\end{theorem}

\begin{proof}[Proof sketch]
\[
  \frac{d}{dt}\mathcal{F}
  =
  \left\langle \frac{\delta \mathcal{F}}{\delta X}, \partial_t X \right\rangle
  +
  \left\langle \frac{\delta \mathcal{F}}{\delta q}, \partial_t q \right\rangle
  =
  -
  \left\langle \frac{\delta \mathcal{F}}{\delta X},
    G_X \frac{\delta \mathcal{F}}{\delta X} \right\rangle
  -
  \left\langle \frac{\delta \mathcal{F}}{\delta q},
    G_q \frac{\delta \mathcal{F}}{\delta q} \right\rangle
  \;\leq\; 0,
\]
since both quadratic forms are nonnegative.
\end{proof}

\paragraph{Interpretation.}
When $\beta = \gamma = 0$, the dynamics reduce to pure RSVP Lyapunov descent.
When the infrastructural state $X$ is fixed, the functional reduces to standard
variational free energy in the style of \citet{friston2010}. In the general
case, both forms of mismatch must contract simultaneously. Thermodynamic
non-closure yields divergence in $E(t)$; epistemic endogeneity yields
convergence to false-belief attractors (Theorem~\ref{thm:epistemic}). Genuine
xylomorphic stability therefore requires descent in both matter and belief.

%% ---------------------------------------------------------------------------
\section{Sycophancy as Curvature Deformation in Epistemic Free Energy}
\label{app:sycophancy}
%% ---------------------------------------------------------------------------

\subsection{Deformation of the Epistemic Landscape}

In the exogenous setting, $\mathcal{F}_{\mathrm{Bayes}}[q;X]$ is minimized when
$q(h)$ matches the posterior induced by the world-conditioned evidence channel,
and its local geometry is truth-tracking: descent reduces inferential mismatch.
Sycophancy alters this geometry by replacing the exogenous joint density
$p(h,o\mid X)$ with an endogenous effective density
$\widetilde{p}_\pi(h,o\mid X,q)$, where $\pi \in [0,1]$ parameterizes
posterior-aligned response bias. The \emph{sycophantically deformed epistemic
free energy} is
\[
  \mathcal{F}_{\mathrm{Bayes}}^{(\pi)}[q;X]
  :=
  \int_{\mathcal{H}}
  q(h)\,\log \frac{q(h)}{\widetilde{p}_\pi(h,o\mid X,q)}\,dh
  =
  \mathcal{F}_{\mathrm{Bayes}}[q;X]
  -
  \pi\,\Psi[q]
  + \mathcal{O}(\pi^2),
\]
where $\Psi[q]$ is a reinforcement functional increasing with posterior
concentration on dominant hypotheses. A canonical choice is
\[
  \Psi[q]
  =
  \tfrac12 \int_{\mathcal{H}}\!\!\int_{\mathcal{H}}
  K(h,h')\,q(h)\,q(h')\,dh\,dh',
\]
with $K(h,h') \geq 0$ a similarity kernel favoring mutually reinforcing
hypotheses.

\begin{definition}[Sycophantic Curvature Deformation]
A sycophantic interaction induces a \emph{negative curvature deformation} if the
Hessian of $\mathcal{F}_{\mathrm{Bayes}}^{(\pi)}$ satisfies
\[
  \delta^2 \mathcal{F}_{\mathrm{Bayes}}^{(\pi)}
  =
  \delta^2 \mathcal{F}_{\mathrm{Bayes}}
  -
  \pi\,\delta^2\Psi
  + \mathcal{O}(\pi^2),
\]
with $\delta^2\Psi$ positive on a nontrivial subspace of perturbations.
\end{definition}

\begin{theorem}[Sycophantic Spurious-Minima Theorem]
\label{thm:sycophantic_minima}
Suppose $\mathcal{F}_{\mathrm{Bayes}}[q;X]$ is locally strictly convex near the
truth-aligned posterior $q^\star$. Let
$\mathcal{F}_{\mathrm{Bayes}}^{(\pi)}[q;X] = \mathcal{F}_{\mathrm{Bayes}}[q;X]
- \pi\,\Psi[q]$ with $\Psi$ twice Fr\'echet differentiable. If, for some
direction $\xi$,
\[
  \pi\,\langle \xi,\delta^2\Psi[q^\star]\xi\rangle
  >
  \langle \xi,\delta^2\mathcal{F}_{\mathrm{Bayes}}[q^\star;X]\xi\rangle,
\]
then $q^\star$ ceases to be a strict local minimum along $\xi$, and the deformed
functional admits either a flat direction or a new off-truth local minimum near
$q^\star$.
\end{theorem}

\begin{proof}[Proof sketch]
The second variation at $q^\star$ gives
$\delta^2 \mathcal{F}_{\mathrm{Bayes}}^{(\pi)}[q^\star](\xi,\xi)
= \delta^2 \mathcal{F}_{\mathrm{Bayes}}[q^\star](\xi,\xi)
- \pi\,\delta^2\Psi[q^\star](\xi,\xi)$.
By hypothesis this becomes nonpositive, so $q^\star$ loses strict local
minimality and nearby spurious minima emerge by standard bifurcation.
\end{proof}

\begin{corollary}
Even if the response channel is restricted to truthful observations, selective
presentation of truths can still reinforce dominant posterior directions, reducing
convexity and generating off-truth minima. Factuality constrains support but does
not restore exogeneity of the evidence channel.
\end{corollary}

\subsection{Critical Sycophancy Threshold and Phase Transition}

Let $q^\star$ be the truth-aligned posterior and define the curvature operators
\[
  H_0 := \delta^2 \mathcal{F}_{\mathrm{Bayes}}[q^\star;X],
  \qquad
  H_\Psi := \delta^2 \Psi[q^\star].
\]

\begin{definition}[Critical Sycophancy Threshold]
\[
  \pi_c
  :=
  \frac{\lambda_{\min}(H_0)}{\lambda_{\max}(H_\Psi)}.
\]
\end{definition}

\begin{theorem}[Sycophancy Phase Transition]
\label{thm:sycophancy_phase}
For the deformed epistemic free energy
$\mathcal{F}_{\mathrm{Bayes}}^{(\pi)} = \mathcal{F}_{\mathrm{Bayes}} - \pi\Psi$:
in the subcritical regime $\pi < \pi_c$, the truth-aligned posterior $q^\star$
remains a strict local minimum; at criticality $\pi = \pi_c$, the landscape
develops a flat direction; in the supercritical regime $\pi > \pi_c$, $q^\star$
ceases to be a local minimum and spurious minima emerge in its neighborhood.
\end{theorem}

\begin{proof}[Proof sketch]
By the Courant--Fischer variational principle,
$\lambda_{\min}(H_0 - \pi H_\Psi) \geq \lambda_{\min}(H_0) -
\pi\,\lambda_{\max}(H_\Psi)$,
which changes sign at $\pi = \pi_c$.
\end{proof}

\paragraph{Regime table.}
The phase transition mirrors the thermodynamic $\lambda$-classification exactly:
\[
\begin{array}{ccc}
\textbf{Thermodynamic} & \quad & \textbf{Epistemic} \\
\hline
\lambda < 1 & \longleftrightarrow & \pi < \pi_c \\
\lambda = 1 & \longleftrightarrow & \pi = \pi_c \\
\lambda > 1 & \longleftrightarrow & \pi > \pi_c
\end{array}
\]
In both cases, a contraction condition defines the boundary between stable and
unstable dynamics. The critical distinction is that epistemic instability does
not require divergence: the system may converge rapidly, but to the wrong
attractor. This is the free-energy signature of delusional spiraling
\citep{chandra2026}: stable descent on a mis-specified landscape.

\paragraph{In practical systems.}
The effective sycophancy parameter $\pi$ may arise from reward shaping,
engagement optimization, or other posterior-reinforcing response policies; the
present analysis treats it as an effective geometric parameter of the evidence
channel rather than a commitment to any single training mechanism. Dual closure
requires both $\lambda < 1$ and $\pi < \pi_c$.

\paragraph{}
A helpful system does not merely descend quickly. It must descend on the right
landscape.

\newpage
%% ---------------------------------------------------------------------------
\bibliographystyle{plainnat}
\begin{thebibliography}{99}

\bibitem[Abramsky and Coecke(2004)]{abramsky2004}
Abramsky, S. and Coecke, B.
\newblock A categorical semantics of quantum protocols.
\newblock In \emph{Proc.\ 19th Annual IEEE Symposium on Logic in Computer
  Science}, pages 415--425, 2004.

\bibitem[Baez and Stay(2010)]{baez2010}
Baez, J.~C. and Stay, M.
\newblock Physics, topology, logic and computation: A rosetta stone.
\newblock In B.~Coecke, editor, \emph{New Structures for Physics}, pages
  95--172. Springer, 2010.

\bibitem[Bennett(1982)]{bennett1982}
Bennett, C.~H.
\newblock The thermodynamics of computation---a review.
\newblock \emph{International Journal of Theoretical Physics},
  21(12):905--940, 1982.

\bibitem[Brookes(1990)]{brookes1990}
Brookes, L.
\newblock The greenhouse effect: The fallacies in the energy efficiency
  solution.
\newblock \emph{Energy Policy}, 18(2):199--201, 1990.

\bibitem[Cheng et al.(2025)]{cheng2025}
Cheng, N., Broadbent, G., and Chappell, W.
\newblock Cognitive loop via in-situ optimization: Self-adaptive reasoning for
  science, 2025.

\bibitem[Daly and Farley(2011)]{daly2011}
Daly, H.~E. and Farley, J.
\newblock \emph{Ecological Economics: Principles and Applications}.
\newblock Island Press, 2nd edition, 2011.

\bibitem[European Union(2023)]{eu2023}
European Union.
\newblock Energy efficiency directive (EU) 2023/1791, 2023.
\newblock \url{https://eur-lex.europa.eu/eli/dir/2023/1791}.

\bibitem[Goodfellow et al.(2016)]{goodfellow2016}
Goodfellow, I., Bengio, Y., and Courville, A.
\newblock \emph{Deep Learning}.
\newblock MIT Press, 2016.

\bibitem[Hilty and Aebischer(2014)]{hilty2014}
Hilty, L.~M. and Aebischer, B.
\newblock ICT for sustainability: An emerging research field.
\newblock \emph{Advances in Intelligent Systems and Computing}, 310:3--36,
  2014.

\bibitem[Horner et al.(2014)]{horner2014}
Horner, J., Harper, P., and Jackson, T.
\newblock The role of ICT in energy use: A review of the literature.
\newblock \emph{Energy Policy}, 70:1--10, 2014.

\bibitem[IEA(2023)]{iea2023}
International Energy Agency.
\newblock Energy efficiency 2023.
\newblock Technical report, IEA, 2023.

\bibitem[Jevons(1865)]{jevons1865}
Jevons, W.~S.
\newblock \emph{The Coal Question}.
\newblock Macmillan and Co., 1865.

\bibitem[Jones et al.(2021)]{jones2021}
Jones, N., Cullen, J., and Cullen, J.~M.
\newblock Rebound effects in energy efficiency investments.
\newblock \emph{Energy Policy}, 152:112182, 2021.

\bibitem[Kauffman(1993)]{kauffman1993}
Kauffman, S.~A.
\newblock \emph{The Origins of Order}.
\newblock Oxford University Press, 1993.

\bibitem[Koomey(2011)]{koomey2011}
Koomey, J.
\newblock Growth in data center electricity use 2005 to 2010.
\newblock Technical report, Analytics Press, 2011.

\bibitem[Lurie(2009)]{lurie2009}
Lurie, J.
\newblock \emph{Higher Topos Theory}.
\newblock Princeton University Press, 2009.

\bibitem[Mac~Lane(1998)]{maclane1998}
Mac~Lane, S.
\newblock \emph{Categories for the Working Mathematician}.
\newblock Springer, 2nd edition, 1998.

\bibitem[Masanet et al.(2020)]{masanet2020}
Masanet, E., Shehabi, A., Lei, N., Smith, S., and Koomey, J.
\newblock Recalibrating global data center energy-use estimates.
\newblock \emph{Science}, 367(6481):984--990, 2020.

\bibitem[Odum(1994)]{odum1994}
Odum, H.~T.
\newblock \emph{Ecological and General Systems}.
\newblock University Press of Colorado, 1994.

\bibitem[O'Dwyer and Malone(2014)]{odwyer2014}
O'Dwyer, K.~J. and Malone, D.
\newblock Bitcoin mining and its energy footprint.
\newblock In \emph{ISSC 2014/CIICT 2014}, 2014.

\bibitem[Rambo and Azevedo(2014)]{rambo2014}
Rambo, C. and Azevedo, I.
\newblock Using data centers for district heating: A review.
\newblock \emph{Energy and Buildings}, 81:123--134, 2014.

\bibitem[Satyanarayanan(2017)]{satyanarayanan2017}
Satyanarayanan, M.
\newblock The emergence of edge computing.
\newblock \emph{Computer}, 50(1):30--39, 2017.

\bibitem[Shams(2025)]{shams2025}
Shams, O.
\newblock The moon should be a computer.
\newblock \emph{Palladium Magazine}, April 2025.

\bibitem[Shehabi et al.(2016)]{shehabi2016}
Shehabi, A. et~al.
\newblock United states data center energy usage report.
\newblock Technical report, Lawrence Berkeley National Laboratory, 2016.

\bibitem[Shulman(2012)]{shulman2012}
Shulman, M.
\newblock Homotopy type theory: The logic of space, 2012.

\bibitem[Sorrell(2009)]{sorrell2009}
Sorrell, S.
\newblock Jevons' paradox revisited.
\newblock \emph{Energy Policy}, 37(4):1456--1469, 2009.

\bibitem[Stockholm Data Parks(2020)]{stockholm2020}
Stockholm Data Parks.
\newblock Sustainable data centers: Using waste heat for district heating.
\newblock \url{https://www.stockholmdataparks.com/}, 2020.

\bibitem[Strubell et al.(2019)]{strubell2019}
Strubell, E., Ganesh, A., and McCallum, A.
\newblock Energy and policy considerations for deep learning in NLP\@.
\newblock \emph{arXiv preprint}, 2019.

\bibitem[Yuan et al.(2024)]{yuan2024}
Yuan, Z., Liu, J., Song, G., and Zhu, T.
\newblock Heat: Satellite's meat is GPU's poison.
\newblock \emph{arXiv preprint}, 2024.

\bibitem[Landauer(1961)]{landauer1961}
Landauer, R.
\newblock Irreversibility and heat generation in the computing process.
\newblock \emph{IBM Journal of Research and Development}, 5(3):183--191, 1961.

\bibitem[Lloyd(2000)]{lloyd2000}
Lloyd, S.
\newblock Ultimate physical limits to computation.
\newblock \emph{Nature}, 406:1047--1054, 2000.

\bibitem[Friston(2010)]{friston2010}
Friston, K.
\newblock The free-energy principle: A unified brain theory?
\newblock \emph{Nature Reviews Neuroscience}, 11(2):127--138, 2010.

\bibitem[Parr et al.(2022)]{parr2022}
Parr, T., Pezzulo, G., and Friston, K.
\newblock \emph{Active Inference: The Free Energy Principle in Mind, Brain, and
  Behavior}.
\newblock MIT Press, 2022.

\bibitem[Jacobson(1995)]{jacobson1995}
Jacobson, T.
\newblock Thermodynamics of spacetime: The Einstein equation of state.
\newblock \emph{Physical Review Letters}, 75(7):1260--1263, 1995.

\bibitem[Verlinde(2011)]{verlinde2011}
Verlinde, E.
\newblock On the origin of gravity and the laws of Newton.
\newblock \emph{Journal of High Energy Physics}, 2011(4):29, 2011.

\bibitem[Risken(1989)]{risken1989}
Risken, H.
\newblock \emph{The Fokker--Planck Equation: Methods of Solution and
  Applications}.
\newblock Springer, 1989.

\bibitem[Villani(2009)]{villani2009}
Villani, C.
\newblock \emph{Optimal Transport: Old and New}.
\newblock Springer, 2009.

\bibitem[Leinster(2014)]{leinster2014}
Leinster, T.
\newblock \emph{Basic Category Theory}.
\newblock Cambridge University Press, 2014.

\bibitem[Spivak(2014)]{spivak2014}
Spivak, D.~I.
\newblock \emph{Category Theory for the Sciences}.
\newblock MIT Press, 2014.

\bibitem[IEA(2024)]{iea2024}
International Energy Agency.
\newblock Electricity 2024: Analysis and forecast to 2026.
\newblock IEA Report, 2024.

\bibitem[NASA(2016)]{nasa2016}
NASA\@.
\newblock Thermal control systems for spacecraft.
\newblock NASA Technical Report, 2016.

\bibitem[Wertz et al.(2011)]{wertz2011}
Wertz, J.~R., Everett, D.~F., and Puschell, J.~J., editors.
\newblock \emph{Space Mission Engineering: The New SMAD}.
\newblock Microcosm Press, 2011.

\bibitem[Garrett(2011)]{garrett2011}
Garrett, T.~J.
\newblock Are there basic physical constraints on future anthropogenic emissions
  of carbon dioxide?
\newblock \emph{Climatic Change}, 104:437--455, 2011.

\bibitem[Georgescu-Roegen(1971)]{georgescu1971}
Georgescu-Roegen, N.
\newblock \emph{The Entropy Law and the Economic Process}.
\newblock Harvard University Press, 1971.

\bibitem[Ayres(1998)]{ayres1998}
Ayres, R.~U.
\newblock Eco-thermodynamics: Economics and the second law.
\newblock \emph{Ecological Economics}, 26(2):189--209, 1998.

\bibitem[Bejan(2016)]{bejan2016}
Bejan, A.
\newblock \emph{The Physics of Life: The Evolution of Everything}.
\newblock St.\ Martin's Press, 2016.

\bibitem[Prigogine(1977)]{prigogine1977}
Prigogine, I.
\newblock \emph{Time, Structure, and Fluctuations}.
\newblock Nobel Lecture, 1977.

\bibitem[Chandra et al.(2026)]{chandra2026}
Chandra, K., Kleiman-Weiner, M., Ragan-Kelley, J., and Tenenbaum, J.~B.
\newblock Sycophantic chatbots cause delusional spiraling, even in ideal
  Bayesians.
\newblock \emph{arXiv preprint arXiv:2602.19141}, 2026.

\bibitem[Cover and Thomas(2006)]{cover2006}
Cover, T.~M. and Thomas, J.~A.
\newblock \emph{Elements of Information Theory}, 2nd edition.
\newblock Wiley-Interscience, 2006.

\bibitem[Gao et al.(2022)]{gao2022}
Gao, P. et~al.
\newblock The carbon footprint of artificial intelligence.
\newblock \emph{Nature Climate Change}, 12:822--829, 2022.

\bibitem[Kolmogorov(1965)]{kolmogorov1965}
Kolmogorov, A.~N.
\newblock Three approaches to the quantitative definition of information.
\newblock \emph{Problems of Information Transmission}, 1(1):1--7, 1965.

\bibitem[Srnicek(2017)]{srnicek2017}
Srnicek, N.
\newblock \emph{Platform Capitalism}.
\newblock Polity Press, 2017.

\bibitem[Zuboff(2019)]{zuboff2019}
Zuboff, S.
\newblock \emph{The Age of Surveillance Capitalism: The Fight for a Human
  Future at the New Frontier of Power}.
\newblock PublicAffairs, 2019.

\end{thebibliography}

\end{document}
